\documentclass[11pt]{article}
\usepackage[T1]{fontenc}
\usepackage{lmodern}
\usepackage[margin=1in]{geometry}
\usepackage{amsmath,amssymb,amsthm}
\usepackage{booktabs}
\numberwithin{equation}{section}
\usepackage{microtype}
\setlength{\emergencystretch}{1.5em}
\usepackage[hidelinks]{hyperref}
\usepackage{xcolor}
\newcommand{\e}{\mathrm{e}}
\newcommand{\ii}{\mathrm{i}}
\newcommand{\dd}{\mathrm{d}}
\newcommand{\Tr}{\operatorname{Tr}}
\newcommand{\Ran}{\operatorname{Ran}}
\newcommand{\id}{\mathbf{1}}
\newcommand{\R}{\mathbb{R}}
\newcommand{\Z}{\mathbb{Z}}
\newcommand{\N}{\mathbb{N}}
\newcommand{\C}{\mathbb{C}}
\newcommand{\Lamb}{\Lambda}
\newcommand{\Lambstar}{\Lambda^{*}}
\newcommand{\ph}{\mathrm{ph}}
\newcommand{\eff}{\mathrm{eff}}
\newcommand{\hop}{\mathrm{hop}}
\newcommand{\pair}{\mathrm{pair}}
\newcommand{\maj}{\mathrm{maj}}
\newcommand{\ODLRO}{\mathrm{ODLRO}}
\newcommand{\PKH}{\mathrm{PKH\text{-}Hol}}
\newcommand{\Gam}{\Gamma}
\newcommand{\fN}{N_{\mathrm{ph}}}
\newcommand{\qedbox}{\ensuremath{\square}}
\DeclareMathOperator{\var}{var}
\newtheorem{theorem}{Theorem}[section]
\newtheorem{lemma}[theorem]{Lemma}
\newtheorem{proposition}[theorem]{Proposition}
\newtheorem{corollary}[theorem]{Corollary}
\theoremstyle{definition}
\newtheorem{definition}[theorem]{Definition}
\newtheorem{remark}[theorem]{Remark}
\allowdisplaybreaks

\title{Positive-temperature pair long-range order in the three-dimensional
$\pi$-flux Penson-Kolb-Hubbard-Holstein model with unbounded phonons}
\author{Davis Rens\\SuperMatics Inc.\\ORCID 0000-0001-5839-5197}
\date{November 21, 2026}

\begin{document}
\maketitle

\begin{abstract}
We extend a recently released theorem on explicit, constructive positive-temperature
onsite pair off-diagonal long-range order in three-dimensional Penson-Kolb-Hubbard
models to the case of on-site Holstein phonons with an unbounded oscillator Fock
space. The pair-hopping interaction remains the only mechanism that can order pairs,
while each site carries a quantum harmonic oscillator that shifts the local charge.
Two parameters control the result: the pair-hopping amplitude $J>0$ and the static
polaronic shift $U_{\mathrm{ph}}=g^{2}/\omega$, where $g$ is the on-site
electron-phonon coupling and $\omega>0$ the phonon frequency. We show three things.
First, reflection positivity, Gaussian domination and the infrared bound survive for
every real $g$, every $\omega>0$, every real $U$ and every real hopping amplitude, on
the full unbounded fermion-phonon Hilbert space: the odd-sublattice particle-hole
transformation used for fermions must be accompanied by the odd-sublattice phonon
parity, and the oscillator cutoff can then be removed by a trace-norm bridge.
Second, the double-commutator sum rule for the bare pair operator acquires an exact
static term $+4U_{\mathrm{ph}}(1-r)$, $r$ being the vacancy density, together with a
density of displaced-oscillator excitations that is bounded by three elementary
estimates. Third, combining the sum rule with a Gibbs variational polaron-dimer bound,
an energy-entropy decomposition at a split phonon frequency, and the vacancy
partition function at zero hopping, yields explicit lower bounds on the staggered pair
order parameter at every temperature $0<T\le T_0$. We list fifteen parameter points
where the lower bound is certified by outward-rounded interval arithmetic, covering
adiabatic to anti-adiabatic phonons, on-site attraction from weak to strong, weak on-site
repulsion shifted into the admissible window by the phonons, and the $\pi$-flux lattice
with $|\kappa|/J\le 0.50$. The bound holds in the thermodynamic limit, with
the oscillator cutoff taken to infinity first at fixed lattice size and then the
lattice size to infinity. No critical temperature, no phase transition, and no
material is claimed.
\end{abstract}

\tableofcontents

\section{Introduction and relation to earlier work}\label{sec:intro}

\subsection{The result being extended}

The Penson-Kolb-Hubbard model \cite{penson1986} places two electrons on a bond with a
pair-hopping amplitude $J$ and charges an on-site interaction $U$. In three
dimensions, with the fermions on the full Fock space and no assumption about
occupation or spin ordering, a recent preprint \cite{rens2026pkodlro} proves explicit
positive-temperature lower bounds on the onsite pair order parameter, in two
parameter regimes: at zero single-electron hopping with pair hopping $J>0$ and
$U\ge-0.5J$, and on the $\pi$-flux lattice with real hopping of amplitude
$|\kappa|\le0.5J$ and $U\ge0$. The method there is a chain of four ingredients:
graded Majorana reflection positivity for the pair interaction in a bipartite
particle-hole basis; the resulting Gaussian domination and Duhamel infrared bound;
the Dyson-Lieb-Simon \cite{dyson1978} and Kennedy-Lieb-Shastry \cite{kennedy1988}
double-commutator sum rule combined with a Falk-Bruch \cite{falk1969} fluctuation
inequality; and two finite-temperature bounds, a Gibbs variational dimer free energy
and an exact operator decomposition of the thermal energy, which close the argument at
a threshold temperature $T_0$ that arises from the positivity of an explicit
strictly convex function (Theorem \ref{thm:cert} below).

The present paper asks whether that chain survives when each lattice site carries a
quantum harmonic oscillator of frequency $\omega$ that couples to the local charge, the
Holstein coupling \cite{holstein1959}. This is not a cosmetic addition. Local phonons
are the standard microscopic origin of an effective attractive interaction between
electrons, and in the strong-coupling limit they generate the polaron and
charge-density-wave physics studied, for the Holstein-Hubbard model, by Miyao
\cite{miyao2016} and, for ground-state spin properties, by Freericks and Lieb
\cite{freericks1995}. Whether such phonons can be combined with a rigorous
positive-temperature pair order bound has, to our knowledge, not been settled. What we
find is that they can, with two structural modifications and one quantitative loss of
reach.

\subsection{What this paper adds}

Three items are new relative to \cite{rens2026pkodlro}.

\begin{enumerate}
\item \emph{Reflection positivity with on-site phonons, at arbitrary coupling.} The
  bipartite particle-hole transformation that turns the pair interaction into a
  positive square also flips the sign of the local charge $n_{x}-1$ on the odd
  sublattice. With a local Holstein coupling $g\,q_{x}(n_{x}-1)$ that transformation
  therefore produces a staggered coupling, which changes sign under the reflection and
  destroys the argument. The cure is to accompany the fermionic transformation by the
  odd-sublattice phonon parity $\prod_{x\in\Lamb_{\text{odd}}}(-1)^{b_{x}^{\dagger}b_{x}}$,
  which flips $q_{x}$ on the odd sublattice and leaves the oscillator number operator
  invariant. The product of the two signs is $+1$, the electron-phonon Hamiltonian is
  then invariant, and the decomposition is again a negative square across the
  reflection plane. This holds for every real $g$, including arbitrarily strong
  coupling, and for every $\omega>0$ (Theorem \ref{thm:rp}).
\item \emph{Unbounded oscillators.} The oscillator space is infinite dimensional, so
  the finite-cutoff reflection-positive matrix inequalities must be shown to pass to
  the limit. We remove the cutoff at fixed lattice size with a trace-norm estimate
  based on a cutoff-uniform lower bound on the local Holstein site Hamiltonian and a
  Dyson-Duhamel expansion in the trace class (Theorem \ref{thm:bridge}). Reflection
  positivity, Gaussian domination, and hence the infrared bound hold on the full
  fermion-phonon space.
\item \emph{Sum rule, energy, and entropy with phonons.} The double commutator of the
  pair operator with the electron-phonon term is exactly on-site and momentum
  independent (Lemma \ref{lem:on-site}); after the polaron completion of the square it
  contributes a static term $+4U_{\mathrm{ph}}(1-r)$ with $U_{\mathrm{ph}}=g^{2}/\omega$
  and $r$ the vacancy density, plus a term controlled by the density $\bar\nu$ of
  displaced-oscillator excitations. We bound $\bar\nu$ by a variational dimer free
  energy with a partial Lang-Firsov \cite{langfirsov1963} displacement, and we control
  the oscillator entropy by splitting the frequency, $\omega=\tfrac12\omega+\tfrac12\omega$,
  so that the thermal oscillator entropy is charged at $\tfrac12\omega$ while the
  energy balance uses the full $\omega$. The result is the bound of Theorem
  \ref{thm:cert} with the certified floors collected in Table \ref{tab:rows}.
\end{enumerate}

\subsection{Order of limits and what a reader should expect}

All statements about the infinite system are obtained in two steps: the oscillator
Fock cutoff is taken to infinity at fixed even lattice size $L$ (so that the model is
the unbounded-oscillator model, and no truncated oscillator algebra is used), and only
then $L\to\infty$ along even side lengths. The order parameter is a $\liminf$ over the
resulting sequences. Consequently every bound in this paper is a statement about
finite-temperature correlations of the untruncated model, not about a truncated
approximation.

The bounds are of the form: for every temperature $0<k_{\mathrm B}T\le T_0$, the
staggered pair order parameter exceeds an explicit positive number. Nothing here
identifies the largest such $T_0$, and nothing here is a statement about a critical
temperature. The fifteen parameter points of Table \ref{tab:rows} are the points at
which the positivity of the explicit function \eqref{eq:final} has been certified with
outward-rounded interval arithmetic.

\subsection{Organisation}

Section \ref{sec:model} fixes the model, the polaron completion of the square, the
order parameter, and the order of limits. Section \ref{sec:rp} contains the
reflection-positivity chain, the Gaussian domination and infrared bounds, the bridge to
unbounded oscillators, and an obstruction showing that a single oscillator shared by a
bond that crosses the reflection plane cannot be handled this way. Section
\ref{sec:sumrule} contains the double-commutator identity, the displaced-oscillator
bounds, and the sum rule. Section \ref{sec:finiteT} contains the variational and the
energy-entropy bounds and the certification scheme. Section \ref{sec:main} states the
main theorem and the table. Section \ref{sec:scope} discusses regimes, the limits of
the method, and lists what is not claimed. Section \ref{sec:supp} describes the
accompanying scripts. The appendices record the numerical verification computations
and the lattice integrals.

\section{Model, conventions, and order of limits}\label{sec:model}

\subsection{Lattice, fields, and Hilbert space}

Let $L\ge2$ be an even integer and let
\begin{equation}
\Lamb:=\left(\Z/L\Z\right)^{3},\qquad |\Lamb|=L^{3},\qquad \Lambstar:=\left(\frac{2\pi}{L}\Z\right)^{3}
\end{equation}
be the periodic cubic lattice and its dual. The reflection planes are the three
coordinate planes $P_{\mu}=\{x_{\mu}=0\}$; each splits
$\Lamb=\Lamb_{-}\sqcup\Lamb_{+}$ into congruent halves and exchanges the even and odd
sublattices,
\begin{equation}
\Lamb_{\text{even}}:=\{x:(-1)^{x_{1}+x_{2}+x_{3}}=+1\},\qquad
\Lamb_{\text{odd}}:=\{x:(-1)^{x_{1}+x_{2}+x_{3}}=-1\},
\end{equation}
and we write $\theta$ for the geometric reflection across a plane $P$. Since $L$ is
even, $\theta$ maps $\Lamb_{\text{even}}$ to $\Lamb_{\text{odd}}$.

Fermions carry spin $\sigma\in\{\uparrow,\downarrow\}$ and are described by
$a_{x,\sigma}$ on the fermion Fock space $\mathcal{F}(\Lamb)$; the on-site algebra is
spanned by the pair operators
\begin{equation}
\Delta_{x}:=a_{x,\downarrow}a_{x,\uparrow},\qquad
\Delta_{x}^{\dagger}=a_{x,\uparrow}^{\dagger}a_{x,\downarrow}^{\dagger},\qquad
\Gam^{(1)}_{x}:=\Delta_{x}^{\dagger}+\Delta_{x},\qquad
\Gam^{(2)}_{x}:=-\ii\big(\Delta_{x}^{\dagger}-\Delta_{x}\big),
\label{eq:gamma-def}
\end{equation}
and by the local charge $n_{x}:=\sum_{\sigma}a^{\dagger}_{x,\sigma}a_{x,\sigma}$. Two
elementary facts will be used repeatedly: on each site
\begin{equation}
\big(\Gam^{(1)}_{x}\big)^{2}=\big(\Gam^{(2)}_{x}\big)^{2}=P_{\mathcal{P}_{x}},\qquad
\big[n_{x}-1,\Gam^{(1)}_{x}\big]=2\ii\Gam^{(2)}_{x},\qquad
\big[n_{x}-1,\Gam^{(2)}_{x}\big]=-2\ii\Gam^{(1)}_{x},
\label{eq:local-ids}
\end{equation}
where $P_{\mathcal{P}_{x}}:=\big(n_{x}-1\big)^{2}$ projects onto the sector
$\mathcal{P}_{x}=\mathrm{span}\{|0\rangle_{x},|\uparrow\downarrow\rangle_{x}\}$ in which
the site is either empty or doubly occupied, and
$P_{\mathcal{V}_{x}}:=\id-P_{\mathcal{P}_{x}}$ projects onto the singly occupied
sector $\mathcal{V}_{x}$, which we call the vacancy sector in the terminology of
\cite{rens2026pkodlro}. Thus $(n_{x}-1)^{2}=P_{\mathcal{P}_{x}}$ and
$P_{\mathcal{P}_{x}}+P_{\mathcal{V}_{x}}=\id$.

Each site carries one oscillator, with $[b_{x},b_{y}^{\dagger}]=\delta_{xy}$ and
\begin{equation}
q_{x}:=\frac{b_{x}+b_{x}^{\dagger}}{\sqrt2},\qquad
p_{x}:=\frac{\ii\,(b_{x}^{\dagger}-b_{x})}{\sqrt2},\qquad [q_{x},p_{y}]=\ii\delta_{xy},
\end{equation}
acting on $L^{2}(\R)$ at each site. The full Hilbert space is
\begin{equation}
\mathcal{H}(\Lamb):=\mathcal{F}(\Lamb)\otimes\bigotimes_{x\in\Lamb}L^{2}(\R),
\end{equation}
and $\mathcal{D}(\Lamb)$ is the subspace where every $b_{x}^{\dagger}b_{x}$ has finite
moments; it is a common core for all operators below. Where the lattice dependence is
inessential we write $\langle\cdot\rangle_{\beta}$ for the Gibbs expectation at inverse
temperature $\beta>0$.

\subsection{The Hamiltonian}

With the pair-hopping amplitude $J>0$, the on-site interaction $U\in\R$, the
single-electron hopping amplitude $\kappa\in\R$, the phonon frequency $\omega>0$ and the
Holstein coupling $g\in\R$, the model is
\begin{equation}
H^{(\Lamb)}_{\PKH}:=H_{\pair}+H_{U}+H_{\hop,\pi}+H_{\ph}+H_{\text{el-ph}},
\label{eq:H}
\end{equation}
with
\begin{align}
H_{\pair}&:=\frac{J}{2}\sum_{x\in\Lamb}\sum_{m=1}^{3}
\Big(\Gam^{(1)}_{x}\Gam^{(1)}_{x+e_{m}}+\Gam^{(2)}_{x}\Gam^{(2)}_{x+e_{m}}\Big),
\label{eq:Hpair}\\
H_{U}&:=\frac{U}{2}\sum_{x\in\Lamb}\big(P_{\mathcal{V}_{x}}-\id\big),
\label{eq:HU}\\
H_{\hop,\pi}&:=\sum_{x\in\Lamb}\sum_{m=1}^{3}\sum_{\sigma\in\{\uparrow,\downarrow\}}
\Big(t^{(\pi)}_{x,x+e_{m}}a_{x,\sigma}^{\dagger}a_{x+e_{m},\sigma}
+\overline{t^{(\pi)}_{x,x+e_{m}}}\,a_{x+e_{m},\sigma}^{\dagger}a_{x,\sigma}\Big),
\label{eq:Hhop}\\
H_{\ph}&:=\omega\sum_{x\in\Lamb}b_{x}^{\dagger}b_{x},\qquad
H_{\text{el-ph}}:=g\sum_{x\in\Lamb}q_{x}\big(n_{x}-1\big).
\label{eq:Hph}
\end{align}
The interaction \eqref{eq:HU} equals $-U\sum_{x}(n_{x,\uparrow}n_{x,\downarrow}-\frac12(n_{x,\uparrow}+n_{x,\downarrow}))$
minus the constant $\frac{U}{2}|\Lamb|$, so no correlation is affected, and as in
\cite{rens2026pkodlro} the coupling $U>0$ is attractive: a doubly occupied site pays
$-\frac{U}{2}$ while a singly occupied one pays $0$, so pair formation lowers the energy,
and $U<0$ is repulsive. The form \eqref{eq:HU} is the one for which the polaron identity
\eqref{eq:polaron} collects the interaction and the electron-phonon term into a single
expression in $U_{\eff}$. The amplitudes in \eqref{eq:Hhop} satisfy $|\kappa|=|t^{(\pi)}_{x,y}|$ and
$t^{(\pi)}_{y,x}=\overline{t^{(\pi)}_{x,y}}$, which is exactly the Hermiticity of
\eqref{eq:Hhop}; in the standard gauge they are purely imaginary,
$t^{(\pi)}_{x,x+e_{m}}=\ii\kappa s^{(m)}_{x}$ with $s^{(m)}_{x}\in\{\pm1\}$, hence
antisymmetric, $t^{(\pi)}_{y,x}=-t^{(\pi)}_{x,y}$, and the product of phases around every
elementary plaquette equals $-1$, that is, the flux is $\pi$ per plaquette, in the sense
of Lieb's flux phase \cite{lieb1994,macris1996}. No explicit gauge is needed below. The
$\kappa=0$ case is called the zero-hopping case; it is the Penson-Kolb-Hubbard model
with Holstein phonons.

Two derived couplings appear throughout:
\begin{equation}
U_{\mathrm{ph}}:=\frac{g^{2}}{\omega}\qquad(\text{static polaronic shift}),
\qquad
U_{\eff}:=U+U_{\mathrm{ph}}\qquad(\text{effective local coupling}).
\label{eq:Uph}
\end{equation}

The hypotheses on the parameters are those of \cite{rens2026pkodlro} transported to the
phonon problem. In the zero-hopping case $\kappa=0$ we require
\begin{equation}
U_{\eff}\ge-\frac{J}{2},
\label{eq:HypA}
\end{equation}
and in the $\pi$-flux case we require
\begin{equation}
U_{\eff}\ge0,\qquad \frac{|\kappa|}{J}\le\frac12 .
\label{eq:HypB}
\end{equation}
These conditions are not needed for reflection positivity, Gaussian domination or the
infrared bound, which hold for all real $U$, $\kappa$ and $g$ (Section \ref{sec:rp}); they
delimit the region in which the quantitative bounds of Sections \ref{sec:sumrule} and
\ref{sec:finiteT} are available. In the zero-hopping case \eqref{eq:HypA} reproduces the
range used in \cite{rens2026pkodlro}, and $U_{\eff}\ge0$ in addition makes the vacancy
partition bound \eqref{eq:rmax} available. When $\kappa\ne0$ the vacancy bound is not
available, so $r_{\max}=1$ is used, and \eqref{eq:HypB} keeps the hopping cost
\eqref{eq:ehop} small compared with the pair energy. Note that \eqref{eq:HypA} refers
to $U_{\eff}$ and not to $U$: the phonons shift the admissible interval of the local
interaction, which is the mechanism behind rows 1 and 2 of Table \ref{tab:rows}.

\subsection{Polaron completion of the square}

Combining the oscillator energy with the Holstein coupling and the local interaction
gives an exact single-site identity, which is the reason the static shift \eqref{eq:Uph}
appears everywhere.

\begin{lemma}[Polaron identity]\label{lem:polaron}
For every real $g$, every $\omega>0$ and every $U$, and with
\begin{equation}
\tilde b_{x}:=b_{x}+\frac{g}{\sqrt2\,\omega}\big(n_{x}-1\big),\qquad
\tilde b_{x}^{\dagger}\tilde b_{x}=b_{x}^{\dagger}b_{x}+\frac{g}{\omega}q_{x}\big(n_{x}-1\big)+\frac{U_{\mathrm{ph}}}{2}P_{\mathcal{P}_{x}},
\label{eq:tildeb}
\end{equation}
one has, on $\mathcal{D}(\Lamb)$,
\begin{equation}
H_{\ph}+H_{\text{el-ph}}+H_{U}
=\omega\sum_{x\in\Lamb}\tilde b_{x}^{\dagger}\tilde b_{x}
+\frac{U_{\eff}}{2}\sum_{x\in\Lamb}\big(P_{\mathcal{V}_{x}}-\id\big),
\label{eq:polaron}
\end{equation}
and, for the single-site term $h_{x}:=\omega b_{x}^{\dagger}b_{x}+g q_{x}(n_{x}-1)+\frac{U}{2}(P_{\mathcal{V}_{x}}-\id)$,
the operator inequality
\begin{equation}
h_{x}\ge\frac{\omega}{2}\,b_{x}^{\dagger}b_{x}-c_{\mathrm{loc}}\,\id,\qquad
c_{\mathrm{loc}}:=\frac{g^{2}}{\omega}+\frac{|U|}{2}.
\label{eq:lower}
\end{equation}
\end{lemma}

\begin{proof}
The first identity is the completion of the square in the oscillator variable:
$\omega b^{\dagger}b+g q(n-1)=\omega\tilde b^{\dagger}\tilde b-\frac{g^{2}}{2\omega}(n-1)^{2}$
with $(n-1)^{2}=P_{\mathcal{P}_{x}}=\id-P_{\mathcal{V}_{x}}$. Hence
$H_{\ph}+H_{\text{el-ph}}=\omega\sum_{x}\tilde b_{x}^{\dagger}\tilde b_{x}+\frac{U_{\mathrm{ph}}}{2}\sum_{x}(P_{\mathcal{V}_{x}}-\id)$,
and adding $H_{U}=\frac{U}{2}\sum_{x}(P_{\mathcal{V}_{x}}-\id)$ gives \eqref{eq:polaron}.

For the inequality, write $\omega b^{\dagger}b=\frac{\omega}{2}b^{\dagger}b+\frac{\omega}{2}b^{\dagger}b$
and complete the square in the second half,
$\frac{\omega}{2}b^{\dagger}b+g q(n-1)=\frac{\omega}{2}\big(b^{\dagger}+\frac{\sqrt2 g}{\omega}(n-1)\big)\big(b+\frac{\sqrt2 g}{\omega}(n-1)\big)-\frac{g^{2}}{\omega}(n-1)^{2}$,
which is a positive square minus a bounded term; together with
$\frac{U}{2}(P_{\mathcal{V}_{x}}-\id)\ge-\frac{|U|}{2}\id$ this gives \eqref{eq:lower}.
\end{proof}

The operator $\tilde b_{x}$ is the annihilation operator of the displaced oscillator,
and $\tilde b_{x}^{\dagger}\tilde b_{x}$ counts displaced (polaronic) excitations. The
identity \eqref{eq:polaron} is the exact statement that phonons and local interaction enter the
Hamiltonian only through $U_{\eff}$ and through the excitation operator $\tilde b_{x}$.

\begin{remark}[Relation to the partially displaced variational state]\label{rem:LF}
A Lang-Firsov transformation \cite{langfirsov1963} with displacement fraction
$\lambda\in[0,1]$,
\begin{equation}
U_{\mathrm{LF}}(\lambda):=\exp\Big(-\frac{\lambda g}{\sqrt2\,\omega}\sum_{x\in\Lamb}\big(b_{x}^{\dagger}-b_{x}\big)\big(n_{x}-1\big)\Big),
\label{eq:LF}
\end{equation}
is unitary, commutes with $n_{x}$, and multiplies a pair-hopping matrix element by the
coherent-state overlap
\begin{equation}
g_{\eff}(\lambda):=J\Big|\langle 0_{x}|\exp\Big(\frac{\sqrt2\lambda g}{\omega}(b_{x}^{\dagger}-b_{x})\Big)|0_{x}\rangle\Big|^{2}
=J\exp\Big(-\frac{2\lambda^{2}g^{2}}{\omega^{2}}\Big),
\label{eq:geff}
\end{equation}
while shifting the local energy to $-\frac{U_{\eff}}{2}+\frac{(1-\lambda)^{2}g^{2}}{2\omega}$
per site. Equation \eqref{eq:geff} is the Franck-Condon suppression of pair hopping by the
phonon cloud and is the origin of the parameter $\lambda$ in Section \ref{sec:finiteT}.
\end{remark}

\subsection{Order parameter and order of limits}

The object bounded below is the staggered, or $\pi$-point, pair correlation. Let
\begin{equation}
\hat\Gam^{(a)}_{\mathbf p}:=|\Lamb|^{-1/2}\sum_{x\in\Lamb}\e^{-\ii\mathbf p\cdot x}\Gam^{(a)}_{x},
\qquad \mathbf p\in\Lambstar ,
\end{equation}
and set $\mathbf Q:=(\pi,\pi,\pi)$.

\begin{definition}[Order parameters]\label{def:order}
For $\beta>0$,
\begin{align}
m_{0}^{2}(T)&:=\liminf_{\substack{L\to\infty\\ L\in2\N}}
\frac{1}{|\Lamb|^{2}}\sum_{x,y\in\Lamb}(-1)^{x_{1}+x_{2}+x_{3}+y_{1}+y_{2}+y_{3}}
\big\langle\Gam^{(1)}_{x}\Gam^{(1)}_{y}\big\rangle^{(\Lamb)}_{\beta},
\label{eq:m0}\\
m^{2}_{\ODLRO}(T)&:=\liminf_{\substack{L\to\infty\\ L\in2\N}}
\frac{1}{|\Lamb|^{2}}\sum_{x,y\in\Lamb}(-1)^{x_{1}+x_{2}+x_{3}+y_{1}+y_{2}+y_{3}}
\big\langle\Delta_{x}^{\dagger}\Delta_{y}\big\rangle^{(\Lamb)}_{\beta}.
\label{eq:modlro}
\end{align}
\end{definition}

Since $\Gam^{(1)}_{x}=\Delta_{x}^{\dagger}+\Delta_{x}$ and the expectations of
$\Delta\Delta$ and $\Delta^{\dagger}\Delta^{\dagger}$ vanish by global $U(1)$ charge conservation,
the two objects differ by a factor of two,
\begin{equation}
m_{0}^{2}(T)=2\,m^{2}_{\ODLRO}(T),
\label{eq:factor2}
\end{equation}
so a positive lower bound on $m_{0}^{2}$ of the form stated in Theorem \ref{thm:cert}
is a lower bound $m^{2}_{\ODLRO}\ge\frac12 m_{0,\mathrm{lo}}^{2}$ on the pair order
parameter itself. We quote both in Table \ref{tab:rows}.

The limits are always taken in the order
\begin{equation}
\fN\to\infty\quad\text{at fixed even }L,\qquad\text{then}\quad L\to\infty\ \text{along even }L .
\label{eq:order}
\end{equation}
The first limit removes the oscillator cutoff, and is the content of Theorem
\ref{thm:bridge}; the second produces the Brillouin-zone integrals of Appendix
\ref{app:integrals}. The order parameter \eqref{eq:modlro} is a two-point correlation of the
$\mathbf Q$ Fourier component, that is, of the staggered pair field; it is the quantity
that is positive in the usual sense of off-diagonal long-range order
\cite{yang1989,lieb1989}. No claim is made about any other Fourier component.

\section{Reflection positivity with unbounded oscillators}\label{sec:rp}

This section contains the first modification of \cite{rens2026pkodlro} required by
phonons: reflection positivity is not automatic, because the particle-hole
transformation that makes the pair interaction a square flips the sign of the local
charge on one sublattice, which produces a staggered Holstein coupling. The cure is a
second unitary, acting on the oscillators.

\subsection{Majorana operators and the graded algebra}

Define, for $x\in\Lamb$ and $\sigma\in\{\uparrow,\downarrow\}$,
\begin{equation}
\begin{aligned}
\xi_{x,\sigma}&:=a_{x,\sigma}+a_{x,\sigma}^{\dagger},\qquad
\eta_{x,\sigma}:=\ii\big(a_{x,\sigma}^{\dagger}-a_{x,\sigma}\big),\\
\{\xi_{x,\sigma},\xi_{y,\tau}\}&=\{\eta_{x,\sigma},\eta_{y,\tau}\}=2\delta_{xy}\delta_{\sigma\tau},\qquad
\{\xi_{x,\sigma},\eta_{y,\tau}\}=0 .
\end{aligned}
\label{eq:maj}
\end{equation}
In this basis the local pair and charge operators are
\begin{equation}
\Gam^{(1)}_{x}=-\frac{\ii}{2}\big(\xi_{x,\uparrow}\eta_{x,\downarrow}+\eta_{x,\uparrow}\xi_{x,\downarrow}\big),\qquad
\Gam^{(2)}_{x}=-\frac{\ii}{2}\big(\xi_{x,\uparrow}\xi_{x,\downarrow}-\eta_{x,\uparrow}\eta_{x,\downarrow}\big),\qquad
n_{x}-1=\frac{\ii}{2}\big(\xi_{x,\uparrow}\eta_{x,\uparrow}+\xi_{x,\downarrow}\eta_{x,\downarrow}\big).
\label{eq:maj-gamma}
\end{equation}
The Majorana monomials of even degree generate the even algebra
$\mathcal{A}_{\mathrm{el}}(\Lamb)$, on which the fermion parity
$\prod_{x}(-1)^{n_{x}}$ acts trivially; the odd monomials form the odd part, and the
$\Z_{2}$ grading is what makes the Majorana reflection positivity of Jaffe and
Pedrocchi \cite{jaffe2015a,jaffe2015b} applicable. The oscillator algebra
$\mathcal{A}_{\ph}(\Lamb)$ is graded even, since it commutes with all Majoranas.

\subsection{The two unitaries}

The first unitary is the bipartite particle-hole transformation that is the basis of
the whole method and that, in the formulation of \cite{koma2022}, is Koma's on-site
gauge transformation. Write $\tilde U_{1}:=\big(U_{1}U_{\ph,\text{odd}}\big)$, where
$U_{1}$ is diagonal in the occupation basis $\{n_{x,\sigma}\}$ and $U_{\ph,\text{odd}}$
sends $\xi_{x,\sigma}\mapsto\xi_{x,\sigma}$, $\eta_{x,\sigma}\mapsto-\eta_{x,\sigma}$ on
the odd sublattice and acts as the identity on the even sublattice. By \eqref{eq:maj-gamma},
\begin{equation}
\Gam^{(1)}_{x}\longmapsto(-1)^{x_{1}+x_{2}+x_{3}}\Gam^{(1)}_{x},\qquad
\Gam^{(2)}_{x}\longmapsto\Gam^{(2)}_{x},\qquad
n_{x}-1\longmapsto(-1)^{x_{1}+x_{2}+x_{3}}\big(n_{x}-1\big),
\label{eq:gammamap}
\end{equation}
so that $H_{U}$, which depends on $P_{\mathcal{P}_{x}}=(n_{x}-1)^{2}$ only, is
invariant, and the pair interaction across a bond becomes a square of differences.
This is exactly the fermionic step in \cite{rens2026pkodlro}.

The problem is the Holstein term. Under $\tilde U_{1}$ alone the coupling becomes
\begin{equation}
g\sum_{x\in\Lamb}q_{x}\big(n_{x}-1\big)
\;\longmapsto\;
g\sum_{x\in\Lamb}(-1)^{x_{1}+x_{2}+x_{3}}q_{x}\big(n_{x}-1\big),
\end{equation}
a staggered coupling which changes sign under the reflection $\theta$ if the oscillator
coordinates are reflected in the natural way $\theta_{\ph}(q_{x})=+q_{\theta(x)}$:
the term is then neither invariant nor a positive square, and the decomposition of the
transformed Hamiltonian into a part supported on $\Lamb_{-}$ plus its reflection is
lost. Reflecting the oscillator antisymmetrically instead,
$\theta_{\ph}(q_{x})=-q_{\theta(x)}$, does not help: for the oscillator observable
$A=q_{x}$ supported on the negative half one then has
$\langle A\,\vartheta(A)\rangle_{\beta}=-\langle q_{x}q_{\theta(x)}\rangle_{\beta}$,
which violates reflection positivity already in the bosonic sector. The computational
check in Appendix \ref{app:checks} reproduces the second failure, with
$\lambda_{\min}=-0.0263317$ in the odd sector for the antisymmetric reflection, on the
small cluster used there.

The second unitary repairs this. On the odd sublattice only, act with the phonon
parity
\begin{equation}
\begin{gathered}
U^{\ph}_{\text{parity,odd}}:=\prod_{x\in\Lamb_{\text{odd}}}(-1)^{b_{x}^{\dagger}b_{x}},
\qquad
U^{\ph}_{\text{parity,odd}}\,b_{x}^{\dagger}b_{x}\,\big(U^{\ph}_{\text{parity,odd}}\big)^{\dagger}=b_{x}^{\dagger}b_{x},\\
U^{\ph}_{\text{parity,odd}}\,q_{x}\,\big(U^{\ph}_{\text{parity,odd}}\big)^{\dagger}
=\begin{cases}+q_{x},&x\in\Lamb_{\text{even}},\\ -q_{x},&x\in\Lamb_{\text{odd}}.\end{cases}
\end{gathered}
\label{eq:parity}
\end{equation}
Under the product unitary
\begin{equation}
\mathcal{U}:=\big(U_{1}U_{\ph,\text{odd}}\big)\otimes U^{\ph}_{\text{parity,odd}}
\label{eq:U}
\end{equation}
the product $q_{x}(n_{x}-1)$ acquires two minus signs on the odd sublattice and none on
the even sublattice, hence
\begin{equation}
\mathcal{U}\,q_{x}\big(n_{x}-1\big)\,\mathcal{U}^{\dagger}=q_{x}\big(n_{x}-1\big)
\quad\text{for all }x\in\Lamb,
\label{eq:invariant}
\end{equation}
and $H_{\ph}+H_{\text{el-ph}}$ and $H_{U}$ are invariant. Note that the parity is
diagonal in the oscillator number basis and therefore commutes with any oscillator
Fock truncation; this is what makes the cutoff argument of Section \ref{sec:bridge}
compatible with \eqref{eq:parity}.

\begin{remark}[Why the parity is not optional]\label{rem:parity}
The two unitaries act on disjoint degrees of freedom and each contributes a sign that
depends on the sublattice. Either one alone produces a staggered, reflection-odd
coupling; only the product is reflection-symmetric. Equivalently, the local Holstein
term is invariant under the two operations only in combination, and the phonon parity
is the oscillator analogue of the odd-sublattice charge flip. The same mechanism
applies to any on-site coupling of the form $q_{x}f(n_{x})$ with $f$ odd about
$n_{x}=1$.
\end{remark}

\subsection{Reflection-positive decomposition}

Let $h=\{h_{m}(x)\in\R\}$ be a collection of real source fields on the nearest-neighbour
bonds and consider the source-dressed Hamiltonian
\begin{equation}
\begin{aligned}
H^{(\Lamb)}_{\PKH}(h):={}&\frac{J}{4}\sum_{x\in\Lamb}\sum_{m=1}^{3}
\Big[\Gam^{(1)}_{x}+\Gam^{(1)}_{x+e_{m}}+h_{m}(x)\id\Big]^{2}
-\frac{J}{4}\sum_{x\in\Lamb}\sum_{m=1}^{3}
\Big[\Gam^{(2)}_{x}-\Gam^{(2)}_{x+e_{m}}\Big]^{2}\\
&+H_{U}+H_{\hop,\pi}+H_{\ph}+H_{\text{el-ph}} .
\end{aligned}
\label{eq:Hh}
\end{equation}
At $h=0$ the on-site squares cancel in \eqref{eq:Hh} and $H^{(\Lamb)}_{\PKH}(0)=H^{(\Lamb)}_{\PKH}$.
The reason for this form is that a source on a bond dresses the pair operator by a
constant, which converts the interaction into a positive square whose minimum is shifted;
the sources are the handle by which Gaussian domination is proved.

\begin{theorem}[Reflection positivity and Gaussian domination with on-site phonons]\label{thm:rp}
Fix an even $L\ge2$, a reflection plane $P$, $\beta>0$, and real parameters
$J>0$, $\omega>0$, $U$, $\kappa$, $g$. Let $\vartheta=\vartheta_{\maj}\otimes\theta_{\ph}$
be the antilinear reflection which acts on Majoranas by
$\vartheta_{\maj}(\xi_{x,\sigma})=\xi_{\theta(x),\sigma}$,
$\vartheta_{\maj}(\eta_{x,\sigma})=-\eta_{\theta(x),\sigma}$ and on oscillators by
reflection of the site labels with complex conjugation in the Fock basis, so that
$\theta_{\ph}(q_{x})=q_{\theta(x)}$. Then:
\begin{enumerate}
\item the transformed Hamiltonian $\widetilde H^{(\Lamb)}(h)=\mathcal{U}H^{(\Lamb)}_{\PKH}(h)\mathcal{U}^{\dagger}$
  has, for every reflection plane and every bond field that vanishes on the bonds
  crossing $P$, the decomposition
  \begin{equation}
  \begin{aligned}
  \widetilde H^{(\Lamb)}(h)={}&\mathcal{A}(h)+\vartheta\big(\mathcal{A}(h)\big)
  -\frac{J}{2}\!\!\sum_{\substack{\langle x,y\rangle\in E_{\text{cross}}\\ x\in\Lamb_{-},\,y=\theta(x)}}\sum_{a=1}^{2}\Gam^{(a)}_{x}\,\vartheta\big(\Gam^{(a)}_{x}\big)\\
  &-\frac{\ii\kappa}{2}\!\!\sum_{\substack{\langle x,y\rangle\in E_{\text{cross}}\\ x\in\Lamb_{-},\,y=\theta(x)}}\sum_{\sigma}\Big(\xi_{x,\sigma}\vartheta(\xi_{x,\sigma})+\eta_{x,\sigma}\vartheta(\eta_{x,\sigma})\Big),
  \end{aligned}
  \label{eq:decomp}
  \end{equation}
  where $\mathcal{A}(h)$ is supported on $\Lamb_{-}$ and contains the internal bonds,
  the local interaction, and the complete on-site oscillator and Holstein terms on
  $\Lamb_{-}$;
\item for every $\beta>0$ the Gibbs state of $\widetilde H^{(\Lamb)}(0)$ is reflection
  positive in the graded sense,
  \begin{equation}
  \Tr\big(A\,\vartheta(A)\,\tilde\rho_{\beta}\big)\ge0\quad(A\in\mathcal{A}(\Lamb_{-})^{\text{even}}),
  \qquad
  \Tr\big(B\,(\ii\vartheta(B))\,\tilde\rho_{\beta}\big)\ge0\quad(B\in\mathcal{A}(\Lamb_{-})^{\text{odd}});
  \label{eq:rp}
  \end{equation}
\item the partition function of the source-dressed Hamiltonian obeys Gaussian domination,
  \begin{equation}
  Z^{(\Lamb)}_{\beta}(h):=\Tr_{\mathcal{H}(\Lamb)}\e^{-\beta H^{(\Lamb)}_{\PKH}(h)}
  \le Z^{(\Lamb)}_{\beta}(0)\qquad\text{for all }h\in\R^{3|\Lamb|}.
  \label{eq:gd}
  \end{equation}
\end{enumerate}
\end{theorem}

\begin{proof}
By \eqref{eq:invariant} the electron-phonon and oscillator parts are invariant under $\mathcal{U}$, and
the local interaction is invariant because it depends on $(n_{x}-1)^{2}$. The
transformed on-site term on $\Lamb_{-}$ is
$\sum_{x\in\Lamb_{-}}\big[\frac{U}{2}(P_{\mathcal{V}_{x}}-\id)+\omega b_{x}^{\dagger}b_{x}+g q_{x}(n_{x}-1)\big]$,
which is (the on-site part of) $\mathcal{A}(h)$; no oscillator operator connects
$\Lamb_{-}$ to $\Lamb_{+}$, because the parity \eqref{eq:parity} is a product of on-site operators.
On a cross-plane bond $\langle x,\theta(x)\rangle$, the pair interaction becomes
\begin{equation}
\frac{J}{4}\Big[\Gam^{(1)}_{x}-\Gam^{(1)}_{\theta(x)}\Big]^{2}-\frac{J}{4}\Big[\Gam^{(2)}_{x}-\Gam^{(2)}_{\theta(x)}\Big]^{2}
=-\frac{J}{2}\sum_{a=1}^{2}\Gam^{(a)}_{x}\vartheta\big(\Gam^{(a)}_{x}\big),
\end{equation}
using $(\Gam^{(a)})^{2}=P_{\mathcal{P}}$ and $\vartheta_{\maj}(\Gam^{(1)})=\Gam^{(1)}$,
$\vartheta_{\maj}(\Gam^{(2)})=-\Gam^{(2)}$, which follow from \eqref{eq:gammamap}; the cross-plane
hopping terms have the form displayed by the same computation as in
\cite{koma2022}. This proves \eqref{eq:decomp}. Since the pair
operators are even and the oscillator algebra is graded even, the decomposition is of
Jaffe-Pedrocchi type, and the reflection-positivity statement of \cite{jaffe2015b} gives
\eqref{eq:rp}. Gaussian domination \eqref{eq:gd} follows from \eqref{eq:rp} by the
Lie-Trotter-Duhamel argument of \cite{dyson1978} applied to each of the three families of coordinate
planes; the phonon coordinates do not enter the source term, so the Duhamel expansion is
the same as in the purely electronic model.
\end{proof}

\begin{remark}[What is and is not used about the phonons]\label{rem:rpreally}
Nothing in Theorem \ref{thm:rp} assumes that $g$ is small, that $\omega$ is large, or
that $U_{\eff}$ has a sign. The phonons enter the reflection-positive structure only
through the parity \eqref{eq:parity}; consequently reflection positivity, Gaussian domination and
the infrared bound below hold for all real $g$ and all $\omega>0$. Where the strength of
the coupling does matter is the quantitative sum rule of Section \ref{sec:sumrule}, not
the positivity structure.
\end{remark}

\subsection{Infrared bound}

Gaussian domination gives a bound on the Duhamel two-point function of the pair field,
in the standard way. Let
\begin{equation}
\big(\hat\Gam^{(a)}_{\mathbf p},\hat\Gam^{(a)}_{-\mathbf p}\big)_{\beta}
:=\frac{1}{Z^{(\Lamb)}_{\beta}}\int_{0}^{1}\Tr\Big(\e^{-(1-s)\beta H}\hat\Gam^{(a)}_{\mathbf p}\e^{-s\beta H}\hat\Gam^{(a)}_{-\mathbf p}\Big)\dd s
\label{eq:duhamel}
\end{equation}
be the Duhamel correlation and $u_{\mathbf p}:=\frac13\sum_{\mu=1}^{3}\cos p_{\mu}$,
$E_{\mathbf p}:=\sum_{\mu=1}^{3}\big(1-\cos p_{\mu}\big)=3\big(1-u_{\mathbf p}\big)$.

\begin{corollary}[Infrared bound with phonons]\label{cor:ir}
For every $\beta>0$, every $\omega>0$, every $J>0$ and all real $U,\kappa,g$, every
$\mathbf p\in\Lambstar\setminus\{\mathbf Q\}$ and $a\in\{1,2\}$,
\begin{equation}
\big(\hat\Gam^{(a)}_{\mathbf p},\hat\Gam^{(a)}_{-\mathbf p}\big)_{\beta}\le\frac{1}{\beta J\,E_{\mathbf p+\mathbf Q}}
=\frac{1}{3\beta J\big(1+u_{\mathbf p}\big)} .
\label{eq:ir}
\end{equation}
\end{corollary}

\begin{proof}
Because the source field $h_{m}(x)$ couples only to the purely fermionic pair
operators $\Gam^{(1)}_{x}+\Gam^{(1)}_{x+e_{m}}$, the second derivative of
$\ln Z^{(\Lamb)}_{\beta}(\varepsilon f)$ at $\varepsilon=0$ is a Duhamel correlation of
pair fields and is non-positive by \eqref{eq:gd}. Taking a staggered source
$\tilde f_{m}(x)=(-1)^{x_{1}+x_{2}+x_{3}}f_{m}(x)$ with
$f_{m}(x)=|\Lamb|^{-1/2}(1-\e^{-\ii k_{m}})\e^{\ii k\cdot x}$ at $k=\mathbf p+\mathbf Q$
and using $E_{\mathbf p+\mathbf Q}=3(1+u_{\mathbf p})$ gives \eqref{eq:ir}. The same bound for
$a=2$ follows from the local pair-phase rotation $\e^{\ii\frac{\pi}{4}n_{x}}$, which
commutes with $n_{x}-1$, $q_{x}$ and $b_{x}^{\dagger}b_{x}$ and therefore leaves the
whole Hamiltonian, including the oscillator terms, invariant.
\end{proof}

\begin{remark}\label{rem:no-phonon-correction}
There is no phonon correction to the constant in \eqref{eq:ir}, for a structural reason: the
source couples to fermionic pair operators only, so the oscillator degrees of freedom
do not appear in the second derivative at all, and the Duhamel expansion treats them as
part of the unperturbed dynamics. This is the sharpest sense in which the Holstein
coupling leaves the infrared structure untouched. The bound itself is the lattice
infrared bound in the form familiar from spin systems \cite{frohlich1976,dyson1978},
here for pair fields rather than spins and with the phonon Fock space left free.
\end{remark}

\subsection{Removing the oscillator cutoff}\label{sec:bridge}

At finite oscillator cutoff $\fN$, replace $b_{x}$, $b_{x}^{\dagger}$ by their
compressions $\Pi_{\fN,x}b_{x}\Pi_{\fN,x}$ on
$\mathcal{H}^{(\fN)}=\mathcal{F}(\Lamb)\otimes\bigotimes_{x}\C^{\fN+1}$ and denote by
$H^{(\Lamb,\fN)}(h)$ the resulting operator. The point of the next theorem is that the
truncated model has the untruncated model as its trace-norm limit at fixed $L$, with
error uniform in $\fN$ for the local part, so that \eqref{eq:rp} and \eqref{eq:gd} survive.

\begin{theorem}[Bridge to unbounded oscillators]\label{thm:bridge}
Fix an even $L\ge2$, $\beta>0$, $\omega>0$ and real $J>0,U,\kappa,g$, and let
$h\in\R^{3|\Lamb|}$. Then as $\fN\to\infty$:
\begin{enumerate}
\item the compressed Gibbs operators converge in trace norm,
  \begin{equation}
  \Big\|\Pi_{\fN}\,\e^{-\beta H^{(\Lamb,\fN)}(h)}\,\Pi_{\fN}-\e^{-\beta H^{(\Lamb)}_{\PKH}(h)}\Big\|_{1}\longrightarrow0;
  \label{eq:tracenorm}
  \end{equation}
\item consequently $Z^{(\Lamb,\fN)}_{\beta}(h)\to Z^{(\Lamb)}_{\beta}(h)$, and the
  truncated Gibbs expectation of every bounded fermion operator, of every polynomial in
  $q_{x},p_{x},b_{x}^{\dagger}b_{x}$ multiplied by a bounded fermion operator, and of
  every pair correlator \eqref{eq:duhamel}, converges to the expectation in the unbounded model;
\item the graded reflection positivity \eqref{eq:rp} and the Gaussian domination inequality
  \eqref{eq:gd} hold for the untruncated Hamiltonian on $\mathcal{H}(\Lamb)$, and the infrared
  bound \eqref{eq:ir} holds as stated.
\end{enumerate}
\end{theorem}

\begin{proof}
Write $H^{(\Lamb)}_{\PKH}(h)=H_{\mathrm{loc}}+W(h)$ with
$H_{\mathrm{loc}}=\sum_{x}h_{x}$ the sum of the commuting single-site operators of
Lemma \ref{lem:polaron} and $W(h)$ purely electronic, bounded with norm
$\|W(h)\|\le M_{\Lamb}(h)$ independent of $\fN$, and commuting with $\Pi_{\fN}$. By
\eqref{eq:lower} each compressed local operator obeys
$h^{(\fN)}_{x}\ge\frac{\omega}{2}\Pi_{\fN,x}b_{x}^{\dagger}b_{x}\Pi_{\fN,x}-c_{\mathrm{loc}}\Pi_{\fN,x}$
with $c_{\mathrm{loc}}$ of \eqref{eq:lower} independent of $\fN$, so the compressed local
semigroup satisfies the cutoff-uniform Schatten bound
\begin{equation}
\Tr\Big(\e^{-\beta H^{(\fN)}_{\mathrm{loc}}}\Big)\le\Big[\frac{4\e^{\beta c_{\mathrm{loc}}}}{1-\e^{-\beta\omega/2}}\Big]^{|\Lamb|}=:Z_{\mathrm{loc},*}(\beta)<\infty .
\label{eq:schatten}
\end{equation}
Since $h^{(\fN)}_{x}\to h_{x}$ in the strong resolvent sense and \eqref{eq:lower} holds uniformly,
$\|\e^{-sH^{(\fN)}_{\mathrm{loc}}}-\e^{-sH_{\mathrm{loc}}}\|_{1}\to0$ for every $s>0$. A
Dyson-Duhamel expansion of $\e^{-\beta H^{(\fN)}}$ in the bounded perturbation
$W_{\fN}=\Pi_{\fN}W(h)\Pi_{\fN}$ about the local generator then has its $k$-th
integrand dominated in trace norm, uniformly in $\fN$, by
$Z_{\mathrm{loc},*}(\beta)\|W(h)\|^{k}$, and each integrand converges in trace norm at
every interior point of the $k$-simplex; dominated convergence in the trace class gives
\eqref{eq:tracenorm}. For the moments in part 2, factor the semigroup at $\beta/2$ and note that
$(b_{x}^{\dagger}b_{x}+\id)^{m}$ times the compressed semigroup stays bounded in
operator norm uniformly in $\fN$ by \eqref{eq:lower}, while it converges strongly; then
Gr\"umm's theorem gives trace-norm convergence after multiplication by the
trace-class $\e^{-\frac{\beta}{2}H}$, and the correlators \eqref{eq:duhamel} follow by a further
application of the same factorization. Part 3 follows by taking $\fN\to\infty$ in
\eqref{eq:rp} and \eqref{eq:gd} at fixed $L$; the left sides converge by part 2, and the closedness of
the cone of reflection-positive functionals, together with the positivity of the
expectations in \eqref{eq:rp}, is preserved. The limit in this last step is legitimate because
the matrix elements in \eqref{eq:rp} involve finitely many oscillator modes at any fixed
fermion number sector.
\end{proof}

\begin{remark}[Order of limits]\label{rem:order}
Theorem \ref{thm:bridge} is the only place where $\fN$ appears, and it is applied before
any lattice limit. All constants in \eqref{eq:gd}, \eqref{eq:ir} and in the sum rule
\eqref{eq:sumrule} are independent of $\fN$; the lattice limit is taken afterwards. This is the order \eqref{eq:order}. We emphasize that
the theorem is a statement about fixed finite $L$; the uniform-in-$L$ statement is the
content of the infrared bound and the sum rule, not of the trace-norm limit.
\end{remark}

\subsection{Bond phonons: an obstruction and where it stops}

The odd-sublattice parity works because the Holstein coupling is on-site. A phonon that
belongs to a single bond does not have this property. The following shows that the
argument of this paper cannot be extended to a single oscillator shared by the two sites
of a bond that crosses the reflection plane, and identifies the class of
phonon-modulated interactions that does remain available.

\begin{proposition}[Bond-coupled phonons]\label{prop:bond}
Consider an extra term $H_{\mathrm{SSH},e}=g_{b}\,q_{e}O_{e}$ with a single oscillator
$q_{e}=(b_{e}+b_{e}^{\dagger})/\sqrt2$ attached to a bond $e=\langle x,y\rangle$ that
crosses $P$, and with $O_{e}$ any non-zero cross-plane bilinear such as
$\Gam^{(1)}_{x}\Gam^{(1)}_{y}+\Gam^{(2)}_{x}\Gam^{(2)}_{y}$, single-electron hopping, or
$(n_{x}-1)(n_{y}-1)$. Then:
\begin{enumerate}
\item if $\theta_{\ph}(q_{e})=-q_{e}$, the bosonic reflection positivity fails for every
  $g_{b}$, since $\langle q_{e}\vartheta(q_{e})\rangle_{\beta}=-\langle q_{e}^{2}\rangle_{\beta}<0$;
\item if $\theta_{\ph}(q_{e})=+q_{e}$, then $q_{e}$ commutes with the fermionic algebra
  and has positive and negative eigenvalues, and in the negative eigensector the
  cross-plane interaction has the wrong sign for $g_{b}>0$, so reflection positivity
  fails for every $g_{b}\ne0$;
\item by contrast, an interaction modulated by a symmetric product of on-site
  oscillator operators,
  \begin{equation}
  H_{\text{pair-mod}}=\frac{J}{2}\sum_{\langle x,y\rangle}f(q_{x})f(q_{y})\Big(\Gam^{(1)}_{x}\Gam^{(1)}_{y}+\Gam^{(2)}_{x}\Gam^{(2)}_{y}\Big),
  \label{eq:modmod}
  \end{equation}
  with $f$ a real polynomial or a bounded Borel function, retains the negative-square
  form across every coordinate plane, with $D_{x,a}=f(\pm q_{x})\Gam^{(a)}_{x}$.
\end{enumerate}
\end{proposition}

\begin{proof}
Part 1 is the observation preceding \eqref{eq:parity}. For part 2, $q_{e}$ acts on the Fock space of
one oscillator and is not a positive operator; compressed to $\fN\ge2$ it has eigenvalues
$\pm\lambda_{k}$ with $\lambda_{k}\ne0$ (at $\fN=2$ the value $\pm1/\sqrt2$), while its
commuting with the fermionic algebra means no fermionic square can repair the sign of
$q_{e}\,C\vartheta(C)$ in the negative eigensector. Part 3: under $\mathcal{U}$ and
$\vartheta$ the cross-plane part of \eqref{eq:modmod} is
$-\frac{J}{2}\sum_{\langle x,\theta(x)\rangle}\sum_{a}D_{x,a}\vartheta(D_{x,a})$ with
$D_{x,a}=f(\pm q_{x})\Gam^{(a)}_{x}$, since $f(q_{x})f(q_{\theta(x)})$ is reflection
symmetric and the two signs arising from $\vartheta$ are those of $\Gam^{(a)}$. This is a
negative square, and the internal bonds lie inside $\mathcal{A}$.
\end{proof}

Appendix \ref{app:checks} records an explicit finite matrix for part 2: on a two-site
cluster with an oscillator at each site truncated to six levels, at $\beta=1.5$, $J=0.2$,
$U=1$, $\omega=1$ and no Holstein coupling, a coupling of strength $1.5$ between the
linear oscillator coordinate spanning the reflection plane and the cross-plane pair
bilinear gives $\lambda_{\min}=-3.943830$ for the symmetric reflection, which is
incompatible with reflection positivity. This is the reason Table \ref{tab:rows}
contains only on-site Holstein couplings.

\section{The double commutator and the sum rule}\label{sec:sumrule}

The second modification is quantitative. The Falk-Bruch fluctuation inequality and the
Kennedy-Lieb-Shastry sum rule require the double commutator of the pair operator with the
Hamiltonian; with phonons this commutator is no longer momentum dependent, but it
acquires a static, temperature-independent piece proportional to $U_{\mathrm{ph}}(1-r)$
and one piece proportional to the displaced-oscillator density.

\subsection{The commutator is on-site and momentum independent}

\begin{lemma}[On-site Holstein double commutator]\label{lem:on-site}
For every $x\in\Lamb$ the identities
\begin{equation}
\big[\Gam^{(1)}_{x},n_{x}-1\big]=-2\ii\Gam^{(2)}_{x},\qquad
\big[\Gam^{(2)}_{x},n_{x}-1\big]=+2\ii\Gam^{(1)}_{x},
\label{eq:dc1}
\end{equation}
\begin{equation}
\big[\Gam^{(a)}_{x},\big[n_{x}-1,\Gam^{(a)}_{x}\big]\big]=-4\big(n_{x}-1\big),\qquad
\sum_{a=1}^{2}\big[\Gam^{(a)}_{x},\big[n_{x}-1,\Gam^{(a)}_{x}\big]\big]=-8\big(n_{x}-1\big)
\label{eq:dc2}
\end{equation}
hold, and consequently
\begin{equation}
\Big[\hat\Gam^{(a)}_{\mathbf p},\big[H_{\ph}+H_{\text{el-ph}},\hat\Gam^{(a)}_{-\mathbf p}\big]\Big]
=-\frac{4g}{|\Lamb|}\sum_{x\in\Lamb}q_{x}\big(n_{x}-1\big)
\qquad\text{for all }\mathbf p\in\Lambstar,\ a\in\{1,2\}.
\label{eq:dc3}
\end{equation}
Averaging over $a$, the total double commutator
$c_{\mathbf p}:=\frac12\sum_{a}\big\langle\big[\hat\Gam^{(a)}_{\mathbf p},\big[H,\hat\Gam^{(a)}_{-\mathbf p}\big]\big]\big\rangle_{\beta}$
equals
\begin{equation}
c_{\mathbf p}=12J\,\mathcal{E}_{1}(\beta)+12J\,\mathcal{E}_{\mathrm{diag}}(\beta)u_{\mathbf p}+c^{(\kappa)}_{\mathbf p}(\beta)
-\frac{4g}{|\Lamb|}\sum_{x\in\Lamb}\big\langle q_{x}(n_{x}-1)\big\rangle_{\beta},
\label{eq:dc4}
\end{equation}
with
\begin{equation}
\begin{aligned}
\mathcal{E}_{1}(\beta)&:=-\frac{1}{6|\Lamb|}\sum_{x\in\Lamb}\sum_{m=1}^{3}\big\langle\Gam^{(1)}_{x}\Gam^{(1)}_{x+e_{m}}+\Gam^{(2)}_{x}\Gam^{(2)}_{x+e_{m}}\big\rangle_{\beta},\\
\mathcal{E}_{\mathrm{diag}}(\beta)&:=\frac{1}{3|\Lamb|}\sum_{\langle x,y\rangle}\big\langle(n_{x}-1)(n_{y}-1)\big\rangle_{\beta}\le1-r(\beta),
\end{aligned}
\label{eq:E1}
\end{equation}
where $r(\beta):=\frac{1}{|\Lamb|}\sum_{x}\langle P_{\mathcal{V}_{x}}\rangle_{\beta}\in[0,1]$
is the vacancy density, and $c^{(\kappa)}_{\mathbf p}$ is the hopping double commutator.
The first sum in \eqref{eq:E1} runs over the three positive directions, so that every bond is
counted once. With this normalization the pair-hopping energy density is exactly
$\frac{1}{|\Lamb|}\langle H_{\pair}\rangle_{\beta}=-3J\mathcal{E}_{1}(\beta)$, because
\eqref{eq:Hpair} gives
$H_{\pair}=\frac{J}{2}\sum_{x}\sum_{m=1}^{3}(\Gam^{(1)}_{x}\Gam^{(1)}_{x+e_{m}}+\Gam^{(2)}_{x}\Gam^{(2)}_{x+e_{m}})$;
this identity is what converts the variational bound on the free energy into a lower
bound on $\mathcal{E}_{1}$ in Section \ref{sec:finiteT}, and the same normalization
enters the sum rule \eqref{eq:sumrule}. The local interaction $H_{U}$ contributes nothing.
\end{lemma}

\begin{proof}
Start with the elementary identities \eqref{eq:dc1} and \eqref{eq:dc2}. Since
$[n_{x}-1,\Delta_{x}^{\dagger}]=2\Delta_{x}^{\dagger}$,
$[n_{x}-1,\Delta_{x}]=-2\Delta_{x}$ and $[\Gam^{(1)}_{x},\Gam^{(2)}_{x}]=2\ii(n_{x}-1)$,
the pair interaction contributes $12J\mathcal{E}_{1}(\beta)+12J\mathcal{E}_{\mathrm{diag}}(\beta)u_{\mathbf p}$
to the double commutator of \eqref{eq:dc4}, with the two coefficients defined there. The oscillator operators
commute with all Majoranas and $[n_{y}-1,\Gam^{(a)}_{x}]=0$ for $y\ne x$, so the only
non-vanishing contribution of the phonon terms to the double commutator is the on-site
one, and summing the Fourier factors $\e^{-\ii\mathbf p\cdot x}\e^{+\ii\mathbf p\cdot x}=1$
gives the oscillator identity \eqref{eq:dc3}. Combining this with the pair and hopping
contributions gives \eqref{eq:dc4}.
\end{proof}

\begin{remark}[The two phonon contributions]\label{rem:twophonon}
Equation \eqref{eq:dc3} says that the oscillator part of the double commutator is proportional to
the thermal expectation of the local electron-phonon energy, with a fixed coefficient and
no momentum dependence: the phonon cloud is local. It is this locality, and not any small
parameter, which keeps the sum rule usable. The two terms in \eqref{eq:dc4} that involve $g$ are
the ones bounded in the next subsection.
\end{remark}

\subsection{Bounds on the displaced-oscillator density}

Define the density of displaced-oscillator excitations
\begin{equation}
\bar\nu(\beta):=\frac{1}{|\Lamb|}\sum_{x\in\Lamb}\big\langle\tilde b_{x}^{\dagger}\tilde b_{x}\big\rangle_{\beta}\ge0 ,
\label{eq:nubar}
\end{equation}
with $\tilde b_{x}$ from \eqref{eq:polaron}. Because the polaron shift is exactly
$q_{x}=\tilde q_{x}-\frac{g}{\omega}(n_{x}-1)$ with
$\tilde q_{x}=(\tilde b_{x}+\tilde b_{x}^{\dagger})/\sqrt2$, the phonon contribution to
\eqref{eq:dc4} splits into a static part and a fluctuation part.

\begin{lemma}[Displaced-oscillator bounds]\label{lem:ph}
For every finite even $L$, every $\beta>0$, every $\omega>0$ and all real $J>0,U,\kappa,g$,
with constants independent of $L$:
\begin{enumerate}
\item exact split,
  \begin{equation}
  -\frac{g}{|\Lamb|}\sum_{x\in\Lamb}\big\langle q_{x}(n_{x}-1)\big\rangle_{\beta}
  =U_{\mathrm{ph}}\big(1-r(\beta)\big)-\frac{g}{|\Lamb|}\sum_{x\in\Lamb}\big\langle\tilde q_{x}(n_{x}-1)\big\rangle_{\beta};
  \label{eq:split}
  \end{equation}
\item variational (Cauchy-Schwarz) bound,
  \begin{equation}
  -\frac{g}{|\Lamb|}\sum_{x\in\Lamb}\big\langle q_{x}(n_{x}-1)\big\rangle_{\beta}
  \le U_{\mathrm{ph}}\big(1-r(\beta)\big)+\sqrt2\,|g|\sqrt{\big(1-r(\beta)\big)\bar\nu(\beta)};
  \label{eq:cs}
  \end{equation}
\item normal-ordered commutator bound,
  \begin{equation}
  -\frac{g}{|\Lamb|}\sum_{x\in\Lamb}\big\langle q_{x}(n_{x}-1)\big\rangle_{\beta}
  \le U_{\mathrm{ph}}\big(1-r(\beta)\big)+\frac{\sqrt2\,|g|}{\omega}C_{\mathrm{comm}}\sqrt{\bar\nu(\beta)}
  +\frac{U_{\mathrm{ph}}}{\omega}C_{\mathrm{comm}}\sqrt{1-r(\beta)},
  \label{eq:comm}
  \end{equation}
  where $C_{\mathrm{comm}}=C_{\mathrm{comm}}(J,\kappa):=12J+2\sqrt6\,|\kappa|$ bounds
  the electronic commutator $K_{x}=[H_{\pair}+H_{\hop,\pi},n_{x}-1]$ in operator norm on
  the full Fock space, uniformly in $U$, $g$, $\omega$ and the volume. Because
  $n_{y}-1$ commutes with every operator at a site $y\ne x$, only the six oriented terms
  of \eqref{eq:Hpair} with one index at $x$ contribute to the pair part. Each carries the
  prefactor $J/2$ and two components; by \eqref{eq:local-ids} and by
  $\|\Gam^{(a)}_{x}\|=1$, the commutator of a component with $n_{x}-1$ has norm $2$, and
  $\|\Gam^{(b)}_{x}\Gam^{(a)}_{y}\|=1$ because the two factors act on different sites. One
  oriented term therefore contributes at most $(J/2)\cdot2\cdot2=2J$, and the pair part
at most $12J$. For the hopping part fix a spin $\sigma\in\{\uparrow,\downarrow\}$ and put
$D_{\sigma}:=\sum_{y\sim x}t^{(\pi)}_{x,y}a_{y,\sigma}$, the sum running over the six nearest
neighbours $y$ of $x$, with $t^{(\pi)}_{y,x}=\overline{t^{(\pi)}_{x,y}}$. Since
$[a_{x,\sigma}^{\dagger}a_{y,\sigma},n_{x}-1]=-a_{x,\sigma}^{\dagger}a_{y,\sigma}$ and
$[a_{y,\sigma}^{\dagger}a_{x,\sigma},n_{x}-1]=a_{y,\sigma}^{\dagger}a_{x,\sigma}$ for $y\ne x$,
the terms of \eqref{eq:Hhop} with one index at $x$ combine as
$[H_{\hop,\pi},n_{x}-1]=X_{\uparrow}+X_{\downarrow}$ with
$X_{\sigma}:=D_{\sigma}^{\dagger}a_{x,\sigma}-a_{x,\sigma}^{\dagger}D_{\sigma}$. The six
operators in $D_{\sigma}$ act on distinct sites and anticommute, and
$|t^{(\pi)}_{x,y}|=|\kappa|$, so $D_{\sigma}^{\dagger}D_{\sigma}+D_{\sigma}D_{\sigma}^{\dagger}=6\kappa^{2}\,\id$.
Because $D_{\sigma}$ anticommutes with $a_{x,\sigma}$ and $a_{x,\sigma}^{\dagger}$, the cross
terms of $X_{\sigma}^{\dagger}X_{\sigma}$ vanish and
$X_{\sigma}^{\dagger}X_{\sigma}=n_{x,\sigma}D_{\sigma}D_{\sigma}^{\dagger}+(\id-n_{x,\sigma})D_{\sigma}^{\dagger}D_{\sigma}\le6\kappa^{2}\,\id$,
where $n_{x,\sigma}$ commutes with $D_{\sigma}D_{\sigma}^{\dagger}$ and $D_{\sigma}^{\dagger}D_{\sigma}$.
Hence $\|X_{\sigma}\|\le\sqrt6\,|\kappa|$, the hopping part is at most $2\sqrt6\,|\kappa|$,
and $\|K_{x}\|\le12J+2\sqrt6\,|\kappa|$ for every real $\kappa$. The minimum in
  \eqref{eq:Bph} uses whichever of the two bounds is smaller for the parameters at hand,
  and $\mathcal{B}_{\ph}$ is increasing in $C_{\mathrm{comm}}$, so the certificate
  evaluates the commutator bound with this constant;
\item the bare second moment,
  \begin{equation}
  \begin{aligned}
  \frac{1}{|\Lamb|}\sum_{x\in\Lamb}\big\langle q_{x}^{2}\big\rangle_{\beta}
  \le{}&2\bar\nu(\beta)+1+\frac{2g^{2}}{\omega^{2}}\big(1-r(\beta)\big)\\
  &+\frac{2|g|}{\omega}\sqrt{2\bar\nu(\beta)\big(1-r(\beta)\big)} .
  \end{aligned}
  \label{eq:q2}
  \end{equation}
\end{enumerate}
Hence the phonon part of \eqref{eq:dc4} is bounded by $4\mathcal{B}_{\ph}(\beta)$ with
\begin{equation}
\begin{aligned}
\mathcal{B}_{\ph}(\beta):={}&U_{\mathrm{ph}}\big(1-r(\beta)\big)\\
&+\min\Big\{\sqrt2\,|g|\sqrt{\big(1-r(\beta)\big)\bar\nu(\beta)},\;
\frac{\sqrt2\,|g|}{\omega}C_{\mathrm{comm}}\sqrt{\bar\nu(\beta)}+\frac{U_{\mathrm{ph}}}{\omega}C_{\mathrm{comm}}\sqrt{1-r(\beta)}\Big\}.
\end{aligned}
\label{eq:Bph}
\end{equation}
\end{lemma}

\begin{proof}
Insert $q_{x}=\tilde q_{x}-\frac{g}{\omega}(n_{x}-1)$ and use $(n_{x}-1)^{2}=P_{\mathcal{P}_{x}}$
with $\langle P_{\mathcal{P}_{x}}\rangle=1-\langle P_{\mathcal{V}_{x}}\rangle$ to get
\eqref{eq:split}. For the Cauchy-Schwarz bound \eqref{eq:cs}, expand
$0\le\big(\tilde b_{x}+\alpha\frac{g}{\sqrt2\omega}(n_{x}-1)\big)^{\dagger}\big(\tilde b_{x}+\alpha\frac{g}{\sqrt2\omega}(n_{x}-1)\big)$
for $\alpha>0$, use $(n_{x}-1)^{2}=P_{\mathcal{P}_{x}}$, take the thermal expectation and
optimize in $\alpha$. For \eqref{eq:comm}, use the equation of motion in the displaced frame,
$[H,\tilde b_{x}^{\dagger}]=\omega\tilde b_{x}^{\dagger}+\frac{g}{\sqrt2\omega}K_{x}$,
and the cyclicity of the trace at finite $L$ to eliminate the time derivative; the
remainder is bounded by the operator Cauchy-Schwarz inequality applied to
$\tilde b_{x}^{\dagger}(n_{x}-1)$. Equation \eqref{eq:q2} follows from
$\langle\tilde q_{x}^{2}+\tilde p_{x}^{2}\rangle_{\beta}=2\bar\nu(\beta)+1$ and the same
Cauchy-Schwarz step. The final display is the definition of $\mathcal{B}_{\ph}$, which
takes the smaller of the two bounds.
\end{proof}

\begin{remark}[Where the static term comes from]\label{rem:static}
The first term of \eqref{eq:split} is the whole story in the adiabatic limit: with the
displaced coordinate $\tilde q_{x}$ set to zero, that is, with the oscillator following
the charge exactly, the left-hand side of \eqref{eq:split} equals
$\frac{g^{2}}{\omega}(1-r)=U_{\mathrm{ph}}\langle(n-1)^{2}\rangle$, which is twice the
static displacement energy $\frac{U_{\mathrm{ph}}}{2}(1-r)$ gained by displacing. This
term carries a plus sign in the sum rule and the Falk-Bruch step below can absorb it only
through the energy range available from the pair interaction; the only way to reduce it
is a negative correlation between the displacement fluctuation $\tilde q_{x}$ and the
charge $n_{x}-1$, which is what the two Cauchy-Schwarz bounds of Lemma \ref{lem:ph}
estimate. In contrast to the fermionic case, the sum rule therefore cannot be improved
by taking $\beta$ large: the static term is temperature independent.
\end{remark}

\subsection{The sum rule}

We need the Kennedy-Lieb-Shastry \cite{kennedy1988} weighted sum rule with the Falk-Bruch
inequality \cite{falk1969}, as used in \cite{rens2026pkodlro}, transported to the
Hamiltonian \eqref{eq:H} with the double commutator \eqref{eq:dc4}.

\begin{theorem}[Sum rule with phonons]\label{thm:sumrule}
Fix $\alpha\in(0,1]$, $\beta>0$, and real $J>0$, $U$, $\kappa$, $g$ with $\omega>0$. Set
\begin{equation}
c\big(r(\beta),\beta\big):=\frac{e_{\hop}\big(r(\beta)\big)}{6J}
+\frac{\mathcal{B}_{\ph}(\beta)}{3J}+\frac{I_{W}}{9\beta^{2}J^{2}},
\label{eq:cdef}
\end{equation}
with $\mathcal{B}_{\ph}$ from \eqref{eq:Bph}, $I_{W}$ the cubic Watson integral of Appendix
\ref{app:integrals}, and
\begin{equation}
e_{\hop}(r):=\min\Big\{2\sqrt6\,|\kappa|\sqrt r,\;2\sqrt3\,|\kappa|I_{\sin}\Big\},
\label{eq:ehop}
\end{equation}
where $I_{\sin}:=\int_{[-\pi,\pi]^{3}}\frac{\dd^{3}p}{(2\pi)^{3}}\big(\sum_{\mu}\sin^{2}p_{\mu}\big)^{1/2}$.
Then in the limit $L\to\infty$ along even $L$,
\begin{equation}
m_{0}^{2}(T)\ge\frac{1}{1+\alpha}\left[\alpha\big(1-r(\beta)\big)+\mathcal{E}_{1}(\beta)
-\sqrt{J_{\mathrm{sq}}(\alpha)\Big(\mathcal{E}_{1}(\beta)+c\big(r(\beta),\beta\big)\Big)}\right],
\label{eq:sumrule}
\end{equation}
with the squared-weight integral
\begin{equation}
J_{\mathrm{sq}}(\alpha):=\int_{[-\pi,\pi]^{3}}\frac{\dd^{3}p}{(2\pi)^{3}}\frac{\big(\alpha-u_{\mathbf p}\big)_{+}^{2}}{1+u_{\mathbf p}},
\qquad u_{\mathbf p}=\frac13\sum_{\mu}\cos p_{\mu},\quad x_{+}:=\max(x,0).
\label{eq:Jsq}
\end{equation}
\end{theorem}

\begin{proof}
Let $\mathfrak g_{\mathbf p}:=\langle\hat\Gam^{(1)}_{\mathbf p}\hat\Gam^{(1)}_{-\mathbf p}\rangle_{\beta}$
and $\mathfrak b_{\mathbf p}$ its Duhamel counterpart. The Falk-Bruch inequality gives
$\mathfrak g_{\mathbf p}\ge\mathfrak b_{\mathbf p}\ge0$ and
$\mathfrak g_{\mathbf p}^{2}\le\mathfrak b_{\mathbf p}^{2}+\frac{\beta}{4}\mathfrak b_{\mathbf p}\mathfrak c_{\mathbf p}$,
where $\mathfrak c_{\mathbf p}:=\frac12\sum_{a}c^{(a)}_{\mathbf p}$. With the infrared
bound \eqref{eq:ir}, $\mathfrak b_{\mathbf p}\le(3\beta J(1+u_{\mathbf p}))^{-1}$, so
\begin{equation}
\big(1+u_{\mathbf p}\big)\mathfrak g_{\mathbf p}^{2}\le\frac{1}{12J}\,\mathfrak c_{\mathbf p}
+\frac{1}{9\beta^{2}J^{2}\big(1+u_{\mathbf p}\big)}
\qquad(\mathbf p\ne\mathbf Q).
\label{eq:fb}
\end{equation}
Summing $\alpha(1-r)+\mathcal{E}_{1}$ over the Brillouin zone and isolating the
$\mathbf Q$ component, where $u_{\mathbf Q}=-1$, gives
\begin{equation}
\alpha\big(1-r\big)+\mathcal{E}_{1}\le\frac{1+\alpha}{|\Lamb|}\mathfrak g_{\mathbf Q}
+\frac{1}{|\Lamb|}\sum_{\mathbf p\ne\mathbf Q}\big(\alpha-u_{\mathbf p}\big)_{+}\mathfrak g_{\mathbf p},
\end{equation}
and the Cauchy-Schwarz inequality with weights $\big(\alpha-u_{\mathbf p}\big)_{+}/\sqrt{1+u_{\mathbf p}}$
and $\sqrt{1+u_{\mathbf p}}$ bounds the second term by
$\sqrt{J_{\mathrm{sq}}(\alpha)}$ times the square root of
$\frac{1}{12J}\frac{1}{|\Lamb|}\sum_{\mathbf p}\mathfrak c_{\mathbf p}+I_{W}/(9\beta^{2}J^{2})$,
where we used that $\frac{1}{|\Lamb|}\sum_{\mathbf p}(1+u_{\mathbf p})^{-1}\to I_{W}$ by the
symmetry $u_{\mathbf p+\mathbf Q}=-u_{\mathbf p}$. It remains to evaluate the momentum
average of \eqref{eq:dc4}. The term $12J\mathcal{E}_{1}$ does not depend on $\mathbf p$ and
survives the average, and the diagonal term $12J\mathcal{E}_{\mathrm{diag}}u_{\mathbf p}$
averages to zero because $\sum_{\mathbf p}u_{\mathbf p}=0$. The hopping double commutator
averages, exactly as in the kinetic-only identity of \cite{rens2026pkodlro}, to $-2$ times
the $\pi$-flux hopping energy density of the same lattice, so after division by $12J$ it
contributes at most $e_{\hop}(r)/(6J)$ because
$\langle H_{\hop,\pi}\rangle/|\Lamb|\ge-e_{\hop}(r)$. The oscillator term is at most
$4\mathcal{B}_{\ph}$ by \eqref{eq:Bph}, which is $\mathcal{B}_{\ph}/(3J)$ after division by
$12J$. The average is therefore at most $\mathcal{E}_{1}+c(r,\beta)$ with $c$ of
\eqref{eq:cdef}, and \eqref{eq:sumrule} follows.
\end{proof}

\begin{remark}[Sign and role of each term]\label{rem:sumrule-parts}
Three features are worth separating. The term $\alpha(1-r)$ is the pair gain: a site that
is empty or doubly occupied is available to the pair field and a singly occupied site is
not, so this term grows with the density $1-r$ of sites outside the vacancy subspace, and
it is the only one in the bracket that does. The term $\mathcal{E}_{1}$ is the pair
energy: it is the quantity that the sum rule converts into the momentum-space
correlation, at the price of the square root. The square root is the fluctuation cost: it
contains the static phonon part $4U_{\mathrm{ph}}(1-r)$ of \eqref{eq:dc4}, divided by the
same $12J$ that multiplies $\mathcal{E}_{1}$ and therefore entering as
$\mathcal{B}_{\ph}/3J$, together with the hopping cost $e_{\hop}/6J$ and the temperature
cost $I_{W}/(9\beta^{2}J^{2})$. The presence of $U_{\mathrm{ph}}$ inside the square root,
with a plus sign, is exactly the quantitative loss of reach produced by the phonons.
\end{remark}

\section{From the sum rule to explicit floors}\label{sec:finiteT}

Three quantities in \eqref{eq:sumrule} are not known explicitly at finite temperature: the pair energy
$\mathcal{E}_{1}(\beta)$, the vacancy density $r(\beta)$, and the displaced-oscillator
density $\bar\nu(\beta)$ that enters the square root. This section bounds them. The pair
energy is bounded from below by a Gibbs variational bound, the vacancy density from above
by a partition-function computation at zero hopping, and $\bar\nu$ from above jointly with
$\mathcal{E}_{1}$ by an exact operator decomposition of the thermal energy. The result is
a bound that involves only the temperature and the parameters, and which can therefore be
certified numerically.

\subsection{Variational polaron-dimer bound}

\begin{lemma}[Gibbs variational bound with partial polaron displacement]\label{lem:var}
Fix $\lambda\in[0,1]$, a field $h\ge0$ and $\beta_0>0$. With $g_{\eff}(\lambda)=J\e^{-2\lambda^{2}g^{2}/\omega^{2}}$
from \eqref{eq:geff}, define the exact four-level dimer data
\begin{equation}
s(h):=\sqrt{\tfrac14+4h^{2}},\qquad e_{\max}(h):=\tfrac12+s(h),\qquad
Z_{4}(b,h):=1+\e^{-b e_{\max}(h)}+\e^{-2b s(h)}+\e^{-b(e_{\max}(h)+1)},
\end{equation}
\begin{equation}
m_{1}(b,h):=\frac{2h}{s(h)}\frac{1-\e^{-2bs(h)}}{Z_{4}(b,h)} .
\end{equation}
Let $f(\beta):=-\frac{1}{\beta|\Lamb|}\ln\Tr\e^{-\beta H^{(\Lamb)}_{\PKH}}$ be the free
energy per site. The trial state constructed in the proof has energy per site
$e_{\mathrm{trial}}(h,\lambda)$ and entropy per site $s_{\mathrm{trial}}(h,\lambda)$,
neither of which depends on $\beta$, and $0\le s_{\mathrm{trial}}(h,\lambda)\le\ln2$.
Then
\begin{enumerate}
\item for every $\beta>0$ the Gibbs variational inequality reads
  \begin{equation}
  f(\beta)\le e_{\mathrm{trial}}(h,\lambda)-\frac{1}{\beta}s_{\mathrm{trial}}(h,\lambda);
  \label{eq:gibbstrial}
  \end{equation}
\item at $\beta=\beta_0$ this is $f(\beta_0)\le-\frac{U_{\eff}}{2}+F_{\mathrm{dimer}}(\beta_0;h,\lambda)$
with
\begin{equation}
\begin{aligned}
F_{\mathrm{dimer}}(\beta_0;h,\lambda):=g_{\eff}(\lambda)\Bigg(&-\frac{e_{\max}(h)}{2}-\frac{\ln Z_{4}\big(\beta_0g_{\eff}(\lambda),h\big)}{2\beta_0g_{\eff}(\lambda)}
+h\,m_{1}\big(\beta_0g_{\eff}(\lambda),h\big)\\
&-\frac54 m_{1}\big(\beta_0g_{\eff}(\lambda),h\big)^{2}\Bigg)
+\frac{(1-\lambda)^{2}g^{2}}{2\omega},
\end{aligned}
\label{eq:Fdimer}
\end{equation}
by $e_{\mathrm{trial}}-s_{\mathrm{trial}}/\beta_0=-\frac{U_{\eff}}{2}+F_{\mathrm{dimer}}$.
\end{enumerate}
Consequently, with the rational grid $\lambda\in\{0,\frac15,\frac25,\frac35,\frac45,1\}$
and $h\in\{0,\frac1{20},\dots,\frac{15}{20}\}$,
\begin{equation}
F_{\mathrm{up}}(\beta_0):=\min_{\lambda,h}F_{\mathrm{dimer}}(\beta_0;h,\lambda)
\ge-\frac{1}{\beta_0|\Lamb|}\ln\Tr\e^{-\beta_0H^{(\Lamb)}_{\PKH}}+\frac{U_{\eff}}{2} .
\label{eq:Fup}
\end{equation}
\end{lemma}

\begin{proof}
Take the trial state
$\rho^{(\lambda,h)}_{\mathrm{trial}}=U_{\mathrm{LF}}(\lambda)\big(\bigotimes_{d}\rho^{\mathrm{el}}_{\mathrm{dimer},d}(h)\otimes\bigotimes_{x}|0_{x}\rangle\langle 0_{x}|\big)U_{\mathrm{LF}}(\lambda)^{\dagger}$,
where $d$ runs over $|\Lamb|/2$ disjoint nearest-neighbour dimers along the first axis,
$\rho^{\mathrm{el}}_{\mathrm{dimer},d}(h)$ is the four-state trial state on the pair
subspace $\mathcal{P}_{x}\otimes\mathcal{P}_{y}$ used in \cite{rens2026pkodlro}, and
$|0_{x}\rangle$ is the oscillator vacuum. Because $\rho^{\mathrm{el}}_{\mathrm{dimer},d}$ is
supported where $n_{x}-1=\pm1$, the displaced vacuum is a tensor product of Glauber
coherent states with finite displacement $\alpha_{x}=-\lambda g\tau_{x}/(\sqrt2\omega)$,
so the trial state is in the domain of $H$ and the Gibbs variational principle applies.
The expectation of $H_{\pair}$ is the dimer pair hopping screened by the overlap
$g_{\eff}(\lambda)/J$ of \eqref{eq:geff}, the expectation of $H_{\hop,\pi}$ vanishes on the pair
subspace, and $H_{\ph}+H_{\text{el-ph}}+H_{U}$ contributes
$-\frac{U_{\eff}}{2}+\frac{(1-\lambda)^{2}g^{2}}{2\omega}$ per site by \eqref{eq:polaron} and the
definition of the trial state. The entropy is that of the electronic dimers, since the
oscillator factor is pure: with $b=\beta_0g_{\eff}(\lambda)$ the dimer factor is the Gibbs
state of a four-level system, whose mean energy is $-\partial_{b}\ln Z_{4}$ and whose
entropy is $S_{4}=\ln Z_{4}-b\,\partial_{b}\ln Z_{4}$, so the entropy per site is
$s_{\mathrm{trial}}=S_{4}/2\in[0,\ln2]$. Evaluating the electronic dimer expectation in
the pair subspace at $\beta_0$, as in \cite{rens2026pkodlro}, then gives
$e_{\mathrm{trial}}-s_{\mathrm{trial}}/\beta_0=-\frac{U_{\eff}}{2}+F_{\mathrm{dimer}}(\beta_0;h,\lambda)$
with the closed form \eqref{eq:Fdimer}. Because the trial state does not depend on $\beta$, the
Gibbs inequality gives \eqref{eq:gibbstrial} at every $\beta>0$. Finally
$U_{\eff}$ enters only through the constants, which gives the bound \eqref{eq:Fup}.
\end{proof}

The Franck-Condon factor of \eqref{eq:geff} is the only place where the variational state feels the
oscillator frequency, and the term $\frac{(1-\lambda)^{2}g^{2}}{2\omega}$ is the price of
displacing by $\lambda$ rather than $1$. For $\lambda=1$ the two combine to give the exact
polaron energy $-\frac{U_{\eff}}{2}$ with an exponentially small pair hopping, and for
$\lambda=0$ one recovers the free dimer with the full $U_{\eff}$; intermediate fractions
interpolate, and the minimum over the grid selects the best of them at each parameter
point.

\subsection{Energy-entropy decomposition with a frequency split}

\begin{lemma}[Energy decomposition and entropy at split frequency]\label{lem:split}
Fix $\varepsilon\in(0,1)$, $\beta_0>0$ and $\beta\ge\beta_0$. With $r(\beta)$ and
$\mathcal{E}_{1}(\beta)$ from \eqref{eq:E1} and $\bar\nu(\beta)$ from \eqref{eq:nubar}, on $\mathcal{H}(\Lamb)$:
\begin{enumerate}
\item the thermal energy per site decomposes as
  \begin{equation}
  \frac{1}{|\Lamb|}\big\langle H^{(\Lamb)}_{\PKH}\big\rangle_{\beta}
  =-3J\,\mathcal{E}_{1}(\beta)+\frac{1}{|\Lamb|}\big\langle H_{\hop,\pi}\big\rangle_{\beta}
  +(1-\varepsilon)\omega\,\bar\nu(\beta)+\varepsilon\omega\,\bar\nu(\beta)+\frac{U_{\eff}}{2}r(\beta)-\frac{U_{\eff}}{2};
  \label{eq:energy}
  \end{equation}
\item the fermion-phonon entropy satisfies
  \begin{equation}
  \frac{1}{\beta|\Lamb|}S(\rho_{\beta})-(1-\varepsilon)\omega\,\bar\nu(\beta)
  \le\frac{1}{\beta}s_{\mathrm{el}}\big(r(\beta)\big)-\frac{1}{\beta_0}\ln\Big(1-\e^{-\beta_0(1-\varepsilon)\omega}\Big),
  \qquad s_{\mathrm{el}}(r):=\ln2+H_{2}(r),
  \label{eq:entropy}
  \end{equation}
  with $H_{2}(r)=-r\ln r-(1-r)\ln(1-r)$ the binary entropy.
\end{enumerate}
\end{lemma}

\begin{proof}
Part 1 is \eqref{eq:polaron} together with the definitions of $\mathcal{E}_{1}$ and $r$; the
oscillator energy is split as $\omega\bar\nu=(1-\varepsilon)\omega\bar\nu+\varepsilon\omega\bar\nu$
without approximation. Part 2 uses subadditivity of the von Neumann entropy with respect
to the fermion and displaced-oscillator factors, so that
$S(\rho_{\beta})\le S(\rho^{\mathrm{el}}_{\beta})+\sum_{x}S(\tilde\rho^{\ph}_{\beta,x})$;
the electronic factor is bounded by $\frac{1}{|\Lamb|}S\le\ln2+H_{2}(r)$, which is the
vacancy-conditioned entropy bound of \cite{rens2026pkodlro}, and the single
oscillator factor by the exact maximum of the entropy at fixed energy
$(1-\varepsilon)\omega$, which is $\ln\big(1-\e^{-\beta_0(1-\varepsilon)\omega}\big)^{-1}$
at the reference temperature; since $\beta\ge\beta_0$ the bound with $\beta_0$ dominates.
\end{proof}

The role of the split is to make the two sides compatible: the entropy is charged at the
smaller frequency $(1-\varepsilon)\omega$, while the energy is measured with the full
$\omega$. The remaining part $\varepsilon\omega\bar\nu$ of the oscillator energy in
\eqref{eq:energy} is not charged against the entropy, and it is this part that gives the
upper bound \eqref{eq:numax} on $\bar\nu$; at $\varepsilon=0$ it vanishes and no such bound
follows. In the certified rows
$\varepsilon=\frac12$.

\subsection{Simultaneous bounds and the admissible domain}

Combining Lemma \ref{lem:var} with Lemma \ref{lem:split} gives, for every
$\beta\ge\beta_0$ and every trial pair $(\lambda,h)$, the following lower bound. The
Gibbs state satisfies the exact identity
$f(\beta)=\frac{1}{|\Lamb|}\langle H\rangle_{\beta}-\frac{1}{\beta|\Lamb|}S(\rho_{\beta})$;
the fixed trial state gives \eqref{eq:gibbstrial}; and \eqref{eq:entropy} replaces
$-\frac{1}{\beta|\Lamb|}S(\rho_{\beta})$ by $-(1-\varepsilon)\omega\bar\nu(\beta)$ minus
$\frac{1}{\beta}s_{\mathrm{el}}(r)$ and minus $s_{\mathrm{osc}}$. Substituting into
\eqref{eq:energy}, bounding $\langle H_{\hop,\pi}\rangle/|\Lamb|\ge-e_{\hop}(r)$, dropping
the non-negative $\varepsilon\omega\bar\nu$, and using \eqref{eq:gibbstrial} gives a lower
bound for $3J\mathcal{E}_{1}$ that carries the entropy combination
$\frac{1}{\beta}\big(s_{\mathrm{trial}}-s_{\mathrm{el}}(r)\big)$. That combination is
non-positive, because $s_{\mathrm{trial}}\le\ln2\le s_{\mathrm{el}}(r)$, so replacing
$1/\beta$ by the larger $1/\beta_0$ lowers it, and the resulting
$s_{\mathrm{trial}}/\beta_0$ combines with the trial energy into
$F_{\mathrm{dimer}}=e_{\mathrm{trial}}+U_{\eff}/2-s_{\mathrm{trial}}/\beta_0$; the minimum
over the trial grid turns $-e_{\mathrm{trial}}+s_{\mathrm{trial}}/\beta_0$ into
$-F_{\mathrm{up}}-U_{\eff}/2$. The result is
\begin{equation}
3J\,\mathcal{E}_{1}(\beta)\ge-F_{\mathrm{up}}(\beta_0)-\frac{1}{\beta_0}s_{\mathrm{el}}\big(r(\beta)\big)-s_{\mathrm{osc}}(\beta_0,\varepsilon)
+\frac{U_{\eff}}{2}r(\beta)-e_{\hop}\big(r(\beta)\big)=:E_{\mathrm{lo}}\big(r(\beta)\big),
\label{eq:Elo}
\end{equation}
\begin{equation}
s_{\mathrm{osc}}(\beta_0,\varepsilon):=-\frac{1}{\beta_0}\ln\Big(1-\e^{-\beta_0(1-\varepsilon)\omega}\Big),
\label{eq:sosc}
\end{equation}
and, simultaneously,
\begin{equation}
\begin{aligned}
\bar\nu(\beta)\le{}&\bar\nu_{\max}\big(r(\beta),3J\mathcal{E}_{1}(\beta)\big)\\
:={}&\frac{1}{\varepsilon\omega}\Big(F_{\mathrm{up}}(\beta_0)+\frac{1}{\beta_0}s_{\mathrm{el}}\big(r(\beta)\big)+s_{\mathrm{osc}}(\beta_0,\varepsilon)
-\frac{U_{\eff}}{2}r(\beta)+e_{\hop}\big(r(\beta)\big)+3J\mathcal{E}_{1}(\beta)\Big)_{+}.
\end{aligned}
\label{eq:numax}
\end{equation}
Two further constraints hold. The star bound
\begin{equation}
3J\mathcal{E}_{1}(\beta)\le\sqrt3\,J\big(1-r(\beta)\big)
\label{eq:star}
\end{equation}
follows by writing $6J\mathcal{E}_{1}(\beta)$ as the lattice average of the seven-site star
operator of \cite{rens2026pkodlro}. That operator vanishes on a vacant site, and on a
doubly occupied one it is twice the $XY$ coupling of a spin $\frac12$ to the total spin
of the six neighbours, which has length at most $3$, so its largest eigenvalue is
$2\sqrt3$ there and $6J\mathcal{E}_{1}\le2\sqrt3\,J(1-r)$. Similarly, when $\kappa=0$ and
$U_{\eff}\ge0$, the vacancy partition bound
\begin{equation}
r(\beta)\le\frac{1}{1+\e^{\beta U_{\eff}/2}}\le\frac{1}{1+\e^{\beta_0U_{\eff}/2}}=:r_{\max}(\beta_0)
\label{eq:rmax}
\end{equation}
follows because for $\kappa=0$ the projector $P_{\mathcal{V}_{x}}$ commutes with every
term of the Hamiltonian, and conditioning on it produces the same oscillator spectrum
shifted by $U_{\eff}$. When $\kappa\ne0$ or $U_{\eff}<0$ we set $r_{\max}(\beta_0):=1$.

Consequently the pair $(r(\beta),\mathcal{E}_{1}(\beta))$ lies, for every
$\beta\ge\beta_0$, in the compact admissible set
\begin{equation}
\mathcal{D}(\beta_0):=\Big\{(r,e)\in[0,r_{\max}(\beta_0)]\times\R_{\ge0}:\;\max\Big(0,\frac{E_{\mathrm{lo}}(r)}{3J}\Big)\le e\le\frac{1-r}{\sqrt3}\Big\},
\label{eq:domain}
\end{equation}
where the pair energy is written in units of $3J$, so that $e\le(1-r)/\sqrt3$ is
\eqref{eq:star} divided by $3J$; this is the variable used in the certification. Every
point of $\mathcal{D}(\beta_0)$ can be inserted into \eqref{eq:sumrule}: the resulting
right-hand side depends on the temperature only through $\beta\ge\beta_0$, and every
temperature-dependent quantity on the right is enclosed by its value at $\beta_0$, in the
direction that weakens the bound. The temperature cost $I_{W}/(9\beta^{2}J^{2})$ is
largest at $\beta_0$, the entropy term $\frac{1}{\beta}s_{\mathrm{el}}(r)$ is largest at
$\beta_0$ because $1/\beta\le1/\beta_0$, and the bound $\bar\nu_{\max}$ of \eqref{eq:numax}
is written with both of them at $\beta_0$; so at a fixed point $(r,e)$ of the admissible
set the phonon term $\mathcal{B}_{\ph}$ and with it $c$ are no smaller at $\beta_0$ than at
any $\beta\ge\beta_0$, because $\mathcal{B}_{\ph}$ increases with $\bar\nu$. Conversely,
$r_{\max}(\beta)=(1+\e^{\beta U_{\eff}/2})^{-1}$ is decreasing in $\beta$, so its value at
$\beta_0$ allows the widest range of $r$. Estimating every term at $\beta_0$ therefore
gives a lower bound on the right-hand side of \eqref{eq:sumrule} at all
$\beta\ge\beta_0$, and it suffices to certify at $\beta_0$.

\subsection{Certification over the domain}

The final bound is the positive-part function
\begin{equation}
m_{0}^{2}(T)\ge\frac{1}{1+\alpha}\left[\alpha\big(1-r\big)+e-\sqrt{J_{\mathrm{sq}}(\alpha)\big(e+c(r,e,\beta_0)\big)}\right]
\qquad\text{for all }(r,e)\in\mathcal{D}(\beta_0),
\label{eq:final}
\end{equation}
where $c(r,e,\beta_0)$ is \eqref{eq:cdef} evaluated at $\beta=\beta_0$, with $\bar\nu$ replaced by
its upper bound $\bar\nu_{\max}(r,3Je)$ of \eqref{eq:numax}; on each rectangle of the cover
below it is replaced by its supremum over the rectangle, a constant. Positivity of the right-hand side of \eqref{eq:final} everywhere on
$\mathcal{D}(\beta_0)$ is the certification condition.

The function
\begin{equation}
f_{r}(e):=e-\sqrt{J_{\mathrm{sq}}(\alpha)\big(e+c\big)}
\label{eq:fr}
\end{equation}
is strictly convex in $e$, because $\partial_{e}^{2}f_{r}=\frac14\sqrt{J_{\mathrm{sq}}(\alpha)}(e+c)^{-3/2}>0$,
with a single stationary point at $e^{*}=\frac14J_{\mathrm{sq}}(\alpha)-c$; its minimum on
a closed interval is therefore attained at the stationary point if it lies inside the
interval, and at an endpoint otherwise. In the variable $r$ the function is smooth, the
entropy $H_{2}$ is concave, and $E_{\mathrm{lo}}(r)$ is obtained from the grid minimum,
so that no monotonicity in $r$ is available; consequently the domain is covered by
closed rectangles rather than by a point grid.

The coverage used in the accompanying script consists of $300$ closed strips
$I_{k}=[r_{k},r_{k+1}]$ that exhaust $[0,r_{\max}(\beta_0)]$, each of them subdivided into
$16$ closed subintervals $J_{k,j}$ of the admissible interval
$[\max(0,E_{\mathrm{lo}}/(3J)),(1-r)/\sqrt3]$ evaluated over $I_{k}$. All endpoints are
outward rounded, so that neighbouring rectangles overlap and the closed rectangles
$I_{k}\times J_{k,j}$ cover $\mathcal{D}(\beta_0)$. On each rectangle the binary entropy
$H_{2}$, the two forms of the hopping cost, the linear vacancy term
$\frac{U_{\eff}}{2}r$ and the bound $(1-r)/\sqrt3$ are enclosed by interval arithmetic.
The function $\bar\nu_{\max}$ of \eqref{eq:numax} is increasing in $e$, so its supremum
over $J_{k,j}$ is attained at the upper endpoint, and both chain estimates for
$\mathcal{B}_{\ph}$, the direct one and the double-commutator one, are evaluated there;
their smaller upper bound is kept. The minimum of $f_{r}$ over $J_{k,j}$ is enclosed by
convexity, using the stationary point when it lies inside the subinterval and the
appropriate endpoint otherwise. The certified floor is the minimum over all rectangles of
the lower endpoint of the enclosure of $\big(\alpha(1-r)+f_{r}(e)\big)/(1+\alpha)$. All
arithmetic is carried out with fifty significant decimal digits and outward rounding, so
that each reported floor is a rigorous lower bound under the stated hypotheses.

\begin{remark}[What the interval arithmetic does and does not cover]\label{rem:ia}
The interval computation covers the two-dimensional admissible set
$\mathcal{D}(\beta_0)$ and the lattice integrals of Appendix \ref{app:integrals}, and it
produces a lower bound that holds simultaneously for all $\beta\ge\beta_0$ and all
$(r,e)$ allowed by the two finite-temperature bounds. It does not certify anything about
the model beyond the hypotheses of Theorem \ref{thm:cert}: in particular a positive floor
at a given parameter point is a statement about the sum rule and the two bounds, and the
absence of a positive floor at another point is not a statement about the model.
\end{remark}

\section{Main theorem and certified parameter points}\label{sec:main}

\subsection{The theorem}

\begin{theorem}[Bounds on the pair order parameter at $T\le T_0$]\label{thm:cert}
Let $J>0$, $\omega>0$ and $U,\kappa,g\in\R$, and let $U_{\mathrm{ph}}=g^{2}/\omega$,
$U_{\eff}=U+U_{\mathrm{ph}}$. Assume either
\begin{enumerate}
\item[(i)] $\kappa=0$ and $U_{\eff}\ge-\frac{J}{2}$ (zero hopping), or
\item[(ii)] $\kappa\ne0$, $U_{\eff}\ge0$ and $|\kappa|/J\le\frac12$ ($\pi$-flux lattice).
\end{enumerate}
Fix $\alpha\in\{\frac12,\frac34,1\}$, $\varepsilon=\frac12$ and a threshold $\beta_0>0$
with $T_0=1/\beta_0$. If the certification condition holds, that is if
\begin{equation}
m_{0,\mathrm{lo}}^{2}:=\min_{\substack{0\le k<N_{r}\\0\le j<N_{E}}}\;\inf\left(\frac{\alpha\big(1-I_{k}\big)+f_{I_{k}}\big(J_{k,j}\big)}{\alpha+1}\right)>0
\label{eq:cond}
\end{equation}
with $f$ from \eqref{eq:fr} and $c$ from \eqref{eq:cdef}, then for every $0<T\le T_0$
\begin{equation}
m_{0}^{2}(T)\ge m_{0,\mathrm{lo}}^{2}>0,\qquad\text{hence}\qquad m^{2}_{\ODLRO}(T)\ge\frac12 m_{0,\mathrm{lo}}^{2}>0,
\label{eq:conclusion}
\end{equation}
in the thermodynamic limit taken as in \eqref{eq:order}.
\end{theorem}

\begin{proof}
By Theorem \ref{thm:bridge} the reflection positivity, Gaussian domination and infrared
bound hold on the untruncated oscillator space at every even $L$; by Corollary
\ref{cor:ir} the infrared bound is \eqref{eq:ir}. Lemma \ref{lem:on-site} and Lemma \ref{lem:ph}
give the double commutator with its phonon parts, and Theorem \ref{thm:sumrule} gives
\eqref{eq:sumrule} at fixed $L$. Lemma \ref{lem:var} and Lemma \ref{lem:split} give the three
finite-temperature constraints \eqref{eq:energy}, \eqref{eq:numax} and the star and
vacancy bounds \eqref{eq:star} and \eqref{eq:rmax}, so the pair
$(r,\mathcal{E}_{1})$ lies in $\mathcal{D}(\beta_0)$ for all $\beta\ge\beta_0$. At a
fixed point of $\mathcal{D}(\beta_0)$ the correction $c$ of \eqref{eq:cdef} is no larger at
$\beta$ than at $\beta_0$: $\mathcal{B}_{\ph}$ of \eqref{eq:Bph} increases with
$\bar\nu$, and $\bar\nu(\beta)$ is bounded by the $\beta$-independent $\bar\nu_{\max}$
of \eqref{eq:numax}, while $I_{W}/(9\beta^{2}J^{2})$ decreases in $\beta$. The right-hand
side of \eqref{eq:sumrule} decreases with $c$, so at $\beta\ge\beta_0$ it is at least its
value at $\beta_0$ at the same point of $\mathcal{D}(\beta_0)$, and hence at least the
infimum of the right-hand side at $\beta_0$ over the larger set $\mathcal{D}(\beta_0)$,
which is what the interval enclosure \eqref{eq:final} computes. Positivity of that infimum gives the bound \eqref{eq:conclusion} at fixed $L$,
uniformly in $L$ because all the constants in \eqref{eq:sumrule} and in the
finite-temperature constraints are independent of $L$. Taking
$L\to\infty$ along even $L$ gives the stated $\liminf$. Finally \eqref{eq:factor2} converts the bound
on $m_{0}^{2}$ into the bound on $m^{2}_{\ODLRO}$.
\end{proof}

\begin{remark}[Reading the theorem]\label{rem:reading}
The theorem is a statement of the form: if the explicit function \eqref{eq:final} is positive at a
parameter point, then the pair order parameter is bounded below there at every
temperature below $T_0$. The hypotheses (i) and (ii) are those of the released theorem
without phonons \cite{rens2026pkodlro}: they concern $U_{\eff}$ rather than $U$ because
of \eqref{eq:polaron}, and this is the sense in which phonons change the admissible parameter region.
No claim is made that \eqref{eq:final} is positive at any point where it has not been certified, and
no claim is made about the largest $T_0$ at which it is positive.
\end{remark}

\subsection{The fifteen certified parameter points}

Table \ref{tab:rows} lists the parameter points at which \eqref{eq:final} has been certified with
outward-rounded interval arithmetic. In each row the pair-hopping amplitude is the unit,
$J=1$, the reported floor is the certified lower bound on $m_{0}^{2}(T)$ at all
$0<T\le T_0$, and the last column is the corresponding bound on the pair order
parameter itself, obtained from \eqref{eq:factor2} by halving and truncating.

\begin{table}[htbp]
\centering
\footnotesize
\setlength{\tabcolsep}{2.4pt}
\begin{tabular}{@{}rlrrrrrrrrr@{}}
\toprule
Row & regime & $U/J$ & $|\kappa|/J$ & $\omega/J$ & $g/J$ & $U_{\mathrm{ph}}/J$ & $U_{\eff}/J$ & $T_0/J$ & $m_{0}^{2}\ge$ & $m^{2}_{\ODLRO}\ge$\\
\midrule
1 & zero hopping, weak repulsion & $-0.50$ & $0$ & $2.0$ & $0.25$ & $0.03125$ & $-0.46875$ & $0.065$ & $8.7232\cdot10^{-3}$ & $4.3616\cdot10^{-3}$\\
2 & zero hopping, weak repulsion & $-0.52$ & $0$ & $50.0$ & $1.60$ & $0.0512$ & $-0.4688$ & $0.050$ & $1.4148\cdot10^{-2}$ & $7.0740\cdot10^{-3}$\\
3 & soft phonon & $0$ & $0$ & $0.01$ & $0.002$ & $0.0004$ & $0.0004$ & $0.004$ & $5.0292\cdot10^{-2}$ & $2.5146\cdot10^{-2}$\\
4 & adiabatic phonon & $0$ & $0$ & $0.10$ & $0.025$ & $0.00625$ & $0.00625$ & $0.025$ & $4.1730\cdot10^{-2}$ & $2.0865\cdot10^{-2}$\\
5 & intermediate phonon & $0$ & $0$ & $1.00$ & $0.18$ & $0.0324$ & $0.0324$ & $0.075$ & $1.9195\cdot10^{-2}$ & $9.5975\cdot10^{-3}$\\
6 & fast phonon & $0$ & $0$ & $5.00$ & $0.70$ & $0.098$ & $0.098$ & $0.060$ & $4.3915\cdot10^{-3}$ & $2.1957\cdot10^{-3}$\\
7 & anti-adiabatic phonon & $0$ & $0$ & $25.0$ & $1.60$ & $0.1024$ & $0.1024$ & $0.060$ & $8.3300\cdot10^{-3}$ & $4.1650\cdot10^{-3}$\\
8 & on-site attraction & $4.0$ & $0$ & $2.00$ & $0.40$ & $0.08$ & $4.08$ & $0.300$ & $2.9544\cdot10^{-2}$ & $1.4772\cdot10^{-2}$\\
9 & strong attraction & $10.0$ & $0$ & $5.00$ & $0.75$ & $0.1125$ & $10.1125$ & $0.350$ & $6.7152\cdot10^{-3}$ & $3.3576\cdot10^{-3}$\\
10 & $\pi$-flux, weak hopping & $0$ & $0.02$ & $1.0$ & $0.15$ & $0.0225$ & $0.0225$ & $0.060$ & $1.1220\cdot10^{-2}$ & $5.6100\cdot10^{-3}$\\
11 & $\pi$-flux, weak hopping & $0$ & $0.05$ & $2.0$ & $0.12$ & $0.0072$ & $0.0072$ & $0.030$ & $5.7017\cdot10^{-4}$ & $2.8508\cdot10^{-4}$\\
12 & $\pi$-flux, attractive & $4.0$ & $0.10$ & $2.0$ & $0.35$ & $0.06125$ & $4.06125$ & $0.120$ & $6.1437\cdot10^{-3}$ & $3.0718\cdot10^{-3}$\\
13 & $\pi$-flux, strong attraction & $10.0$ & $0.20$ & $5.0$ & $0.60$ & $0.072$ & $10.072$ & $0.160$ & $5.5618\cdot10^{-3}$ & $2.7809\cdot10^{-3}$\\
14 & $\pi$-flux, strong hopping & $12.0$ & $0.35$ & $10.0$ & $0.35$ & $0.01225$ & $12.01225$ & $0.060$ & $6.8403\cdot10^{-3}$ & $3.4201\cdot10^{-3}$\\
15 & $\pi$-flux, max. hopping & $20.0$ & $0.50$ & $10.0$ & $0.45$ & $0.02025$ & $20.02025$ & $0.080$ & $5.5839\cdot10^{-4}$ & $2.7919\cdot10^{-4}$\\
\bottomrule
\end{tabular}
\caption{Certified rows. Every entry is a lower bound for all $0<T\le T_0$, in the
thermodynamic limit taken as \eqref{eq:order}, at $\varepsilon=\frac12$ and with the weight
$\alpha=\frac12$ except row 8, where $\alpha=\frac34$. The pair-hopping amplitude is the
unit, so all entries are dimensionless multiples of $J$ powers as indicated. The
$U_{\mathrm{ph}}$ and $U_{\eff}$ columns follow from \eqref{eq:Uph}; the last column is half the
preceding floor, truncated, and bounds $m^{2}_{\ODLRO}$ by \eqref{eq:factor2}. Rows 1 and 2 lie in the
zero-hopping regime (i) with $U_{\eff}\ge-J/2$; rows 3 to 9 are also zero hopping; rows 10
to 15 satisfy $|\kappa|/J\le\frac12$ and $U_{\eff}\ge0$ and lie in regime (ii).}
\label{tab:rows}
\end{table}

\begin{remark}[On the size of the floors]\label{rem:floors}
The floors are not small in the sense of numerical resolution; they are small in the sense
of being rigorous lower bounds on a condensate fraction, and they range over two orders of
magnitude across the table. The largest floors occur where the vacancy term is favourable
and the threshold temperature is small (rows 3 to 5); the smallest occur where either the
hopping cost is large (rows 11 and 15) or the vacancy sector is strongly suppressed by
attraction (row 9). In all cases the bound is positive, which is the content of the
theorem.
\end{remark}

\section{Regimes, limits of the method, and what is not claimed}\label{sec:scope}

\subsection{The phonon frequency window}

Rows 3 to 7 sweep the phonon frequency from $\omega=0.01J$ to $\omega=25J$ at weak
Holstein coupling, that is from the adiabatic to the anti-adiabatic regime, and all are
certified. This is a structural feature of the method rather than an accident of the
chosen parameters: the oscillator frequency enters the certified bound only through the
static shift $U_{\mathrm{ph}}=g^{2}/\omega$ and through the oscillator entropy bound
$s_{\mathrm{osc}}(\beta_0,\varepsilon)=-\beta_0^{-1}\ln\big(1-\e^{-\beta_0(1-\varepsilon)\omega}\big)$
of \eqref{eq:sosc}, which at $\varepsilon=\frac12$ is
$-\beta_0^{-1}\ln(1-\e^{-\beta_0\omega/2})$. That term is at most $T_0$ once
$\beta_0\omega\gtrsim1$ and decays exponentially for larger $\beta_0\omega$, while as
$\omega\to0$ it grows only logarithmically,
$s_{\mathrm{osc}}\sim\beta_0^{-1}\ln\big(2/(\beta_0\omega)\big)$. It is not bounded
uniformly as $\omega\to0$ at fixed $T_0$: the logarithm diverges, slowly, so the strict
adiabatic limit is reached by lowering $T_0$ together with $\omega$, which is how rows 3
and 4 are certified. There is no frequency threshold in the statement of the theorem, and
the certified floors do vary with $\omega$, because both the static shift, which grows
like $1/\omega$, and the entropy term do.

\subsection{Phonon-mediated attraction and the boundary of the zero-hopping regime}

The zero-hopping regime requires $U_{\eff}=U+g^{2}/\omega\ge-J/2$, whereas the same
condition reads $U\ge-J/2$ without phonons. Because $U_{\mathrm{ph}}>0$, the admissible
interval of the local interaction is shifted upwards by $g^{2}/\omega$: a model with
$U=-0.52J$, which is outside the interval for the purely electronic theorem, is inside it
once $g^{2}/\omega=0.0512J$. Rows 1 and 2 are two such points, and they are the ones in
which the phonons change the qualitative conclusion rather than the numerical size of the
floor. The mechanism is exactly \eqref{eq:polaron}: the static shift is an energy that the oscillator
can supply, and the pair term only needs the total local coupling to stay above
$-J/2$. This is the precise sense in which on-site phonons act as an attractive
interaction here, and the reason the model is called a Holstein extension rather than a
perturbation of the purely electronic model.

The phonon shift is not, however, a substitute for the pair hopping. If $J=0$ the pair
interaction vanishes and there is no order parameter to bound: the pair energy
$\mathcal{E}_{1}$ of \eqref{eq:E1} is zero and the whole argument is empty. The attraction generated by the phonons lowers the cost of doubly
occupied sites, but the mechanism that orders pairs is the Penson-Kolb term
\cite{penson1986}, as in the released theorem.

\subsection{Strong coupling to the phonons, and where the bound degrades}

The bound degrades as the coupling to the phonons grows, and the reason is visible in
\eqref{eq:split}: the phonon part $\mathcal{B}_{\ph}/3J$ of the difference enters the
square root of \eqref{eq:sumrule} with a plus sign and consumes the energy that the star
bound \eqref{eq:star} supplies. How far the argument reaches is a property of the
enumeration rather than of the model, and it can be measured on the same terms as Table
\ref{tab:rows}. For $U=0$, $\kappa=0$, $\alpha=\frac12$ and $T_0=0.06J$, with the
$300\times16$ cover used in the table, the certified floor is positive only for
$g/J\le0.3169$, $0.4568$ and $0.7499$ at $\omega/J=1$, $2$ and $5$, that is for
$U_{\mathrm{ph}}/J\le0.1004$, $0.1043$ and $0.1125$; at the next coupling in the bisection
the method returns no floor. Beyond them the bare-pair sum rule gives nothing and this
route to positivity is closed.

Where the loss comes from is not the static term in this range. Entering as
$4\mathcal{B}_{\ph}/12J=\mathcal{B}_{\ph}/3J$, the two phonon pieces of
$\mathcal{B}_{\ph}$ contribute to the $c$ of \eqref{eq:cdef} the amounts
$U_{\mathrm{ph}}(1-r)/3J\approx0.0271$ and $\sqrt2\,g\sqrt{(1-r)\bar\nu}/3J\approx0.0447$
at $\omega=J$, at the minimising rectangle $r\approx0.190$, with the temperature term
$6.1\times10^{-4}$; the displaced-oscillator piece is the larger of the two and is the one
that introduces the loss as the coupling is turned on. The same decomposition holds at the
other two frequencies, with the same ordering and a minimising rectangle moving to smaller
$r$: $0.0283$ against $0.0467$ at $\omega=2J$, at $r\approx0.186$, and $0.0308$ against
$0.0502$ at $\omega=5J$, at $r\approx0.178$. The upper end of the enclosed $c$ of
\eqref{eq:cdef} at these rectangles is $0.0724$, $0.0756$ and $0.0817$ at
$\omega/J=1$, $2$ and $5$. The script \texttt{probe\_static.py} repeats
the certificate of \eqref{eq:final} rectangle by rectangle and prints these numbers
together with the rectangle at which they are attained.

Two rows of Table \ref{tab:rows} are close to that boundary in this sense. Row 6, at
$\omega=5J$ and $U_{\mathrm{ph}}/J=0.0980$, lies at $93.3\%$ of the boundary coupling for
its own parameters, and row 7, at $\omega=25J$ and $U_{\mathrm{ph}}/J=0.1024$, lies at
$87.0\%$. Their windows are therefore close to the degradation in the coupling, although
their thresholds are far from it in temperature. The remaining rows are certified at
smaller $U_{\mathrm{ph}}/J$, or at $|U|=10$ where the pair interaction supplies more
energy, or with $\kappa\ne0$; the location of the boundary is not worked out for those.

In the regime $U_{\mathrm{ph}}\gtrsim J$ the physically relevant object is not the bare
pair but the dressed one: an electron that carries its own phonon cloud, described by the
Lang-Firsov transformation \cite{langfirsov1963} and its generalizations. The variational
state of Lemma \ref{lem:var} already contains a partial version of that dressing through
$g_{\eff}(\lambda)$, but the sum rule is written for the bare pair operator
$\Delta_{x}=a_{x,\downarrow}a_{x,\uparrow}$, and the Franck-Condon physics of a dressed pair,
together with the competing charge-density-wave order that dominates that regime (see
\cite{miyao2016} for the Holstein-Hubbard case), is not treated here.

\subsection{What the paper does not claim}

\begin{itemize}
\item \emph{A critical temperature.} The threshold $T_0$ in Theorem \ref{thm:cert} is a
  value at which the explicit function \eqref{eq:final} is positive. The bound is a statement at
  $T\le T_0$, and no monotonicity in $T$ of any correlation, and no transition, follows.
\item \emph{The bare-pair sum rule beyond a finite coupling.} As measured in the previous
  subsection, the certified floor is positive only for $U_{\mathrm{ph}}/J\le0.1004$,
  $0.1043$ and $0.1125$ at $\omega/J=1$, $2$ and $5$ for $U=0$ and $\kappa=0$, and it
  fails above those couplings. Beyond the certified rows, no positivity is claimed, and
  the limit quoted is a limit of this enumeration: a sharper bound on
  $\bar\nu(\beta)$, or a sum rule written for a dressed operator, would move it.
\item \emph{Polaron and charge-density-wave physics.} The regime in which the phonon cloud
  is strongly bound to the electron, and the competing charge order that appears there,
  are not treated. Nothing here excludes, or confirms, that order.
\item \emph{Dressed-pair off-diagonal long-range order.} Whether the polaron-dressed pair
  operator has long-range order at positive temperature is an open problem. The sum rule
  proved here concerns the bare pair operator, and no statement about a dressed operator
  is made.
\item \emph{Lattices other than the $\pi$-flux one.} Only the $\pi$-flux lattice with
  $|\kappa|/J\le\frac12$ is covered in the hopping regime, together with the zero-hopping
  case. The zero-flux lattice with $\kappa\ne0$ stays open: the reflection-positive
  decomposition \eqref{eq:decomp} fails there, as the check in Appendix \ref{app:checks}
  shows, where reversing the sign of the $\pi$-flux hopping on one bond, which leaves a
  zero-flux plaquette, drives the odd-sector Gram matrix negative; and no other route to
  a spatial pair infrared bound is available to us.
\item \emph{Other parameter points.} The fifteen rows of Table \ref{tab:rows} are the
  certified parameter points. A parameter point absent from the table is not a point at
  which the order parameter vanishes; it is a point at which this method has not been
  certified, either because some hypothesis of Theorem \ref{thm:cert} fails there or
  because the certified function \eqref{eq:final} is not positive under the interval enclosure.
\item \emph{Materials and superconductivity.} No material, and no identification of the
  model with any superconductor, is claimed. The order parameter bounded here is the
  staggered pair correlation of a lattice model, and the interaction is a pair hopping,
  not an electronic Coulomb interaction.
\end{itemize}

\subsection{Open problems}

Three developments would extend the reach of the argument, in increasing order of
difficulty. The first is the zero-flux case, where the reflection-positive decomposition
fails and any progress must come from a different positivity structure. The second is the
strong-coupling regime, where the sum rule would have to be written for a dressed pair
operator whose commutator with the Hamiltonian is not on-site. The third is the
combination of the two, where both obstructions are present and where the physics is that
of the Holstein-Hubbard model at strong coupling. In each case the structural part of this
paper, namely the reflection-positive decomposition with the odd-sublattice parity and the
trace-norm bridge to unbounded oscillators, is available without change; the quantitative
part would have to be redone.

\section{Supporting material and disclosures}\label{sec:supp}

This preprint has not been peer reviewed.

The accompanying replay supplement contains five Python scripts. The script
\texttt{replay.py} reproduces the finite checks: interval enclosures of the lattice
integrals of Appendix \ref{app:integrals}, exact finite-matrix operator identities, the
reflection-positive Gram matrices at oscillator cutoffs $\fN=4,6,8,10$, the two
obstruction checks of Section \ref{sec:rp}, and the fifteen certified rows of Table
\ref{tab:rows} with their interval enclosures. The script \texttt{probe\_xover.py} imports
that certificate and locates the positivity boundary of Section \ref{sec:scope} by
bisection in the Holstein coupling, at the three frequencies quoted there and for rows 6
and 7 of the table. The script \texttt{probe\_static.py} repeats that certificate rectangle
by rectangle and prints the decomposition of the cost quoted in Section \ref{sec:scope},
at the boundary coupling and at the minimising rectangle. The script \texttt{mutants.py}
applies seven semantic modifications of
the certified computation, each of which must be rejected. The script \texttt{run\_all.py}
runs the suites and reports their exit codes. The digest list \texttt{SHA256SUMS} contains
the SHA-256 digest of every file in the supplement.

Passing these checks establishes no analytic statement: they verify arithmetic, matrix
identities on small systems, and the closure of the certified enclosures. The proofs of
Sections \ref{sec:model} to \ref{sec:main} are independent of the scripts.

\paragraph{Use of software and language models.}
AI agents generated proofs, certificates, code, and manuscript drafts and edits under
D.~Rens' direction and review. Numerical enclosures used \texttt{mpmath} interval
arithmetic, and the linear algebra checks used \texttt{numpy} and \texttt{scipy}. The
analytic statements are presented with complete arguments, but no hand verification of
every algebraic step was performed, and readers relying on any particular bound should
check it against the proof as a whole.

\paragraph{Data availability.}
The released preprint \cite{rens2026pkodlro} and its replay supplement
\cite{rens2026pkodlrosupp} contain the corresponding results for the model without
phonons, and the present paper is written so that the two can be read in parallel. The
replay scripts described above accompany this paper. This preprint is archived at
\href{https://doi.org/10.5281/zenodo.24433375}{\nolinkurl{doi:10.5281/zenodo.24433375}}, its
replay supplement at
\href{https://doi.org/10.5281/zenodo.24433377}{\nolinkurl{doi:10.5281/zenodo.24433377}}, and
the exact released source tree, together with the released PDF, at
\href{https://github.com/SuperMatics-Inc/penson-kolb-holstein-pair-odlro}{\nolinkurl{github.com/SuperMatics-Inc/penson-kolb-holstein-pair-odlro}}.

\appendix

\section{Numerical verification computations}\label{app:checks}

This appendix records what the accompanying scripts compute, with the values they produce.
All of it is finite-dimensional or arithmetic in character; none of it is used in the
proof of any analytic statement.

\paragraph{Operator identities.}
On the four-dimensional on-site fermion space, the maximum absolute deviation from the
identities \eqref{eq:dc1} and \eqref{eq:dc2}, and from $(n_{x}-1)^{2}=P_{\mathcal{P}_{x}}$ together with
$(n_{x}-1)P_{\mathcal{P}_{x}}=(n_{x}-1)$, are all $0.00\times10^{0}$, that is zero to
machine precision.

\paragraph{Reflection-positive Gram matrices at finite cutoff.}
On the two-site cluster with $\pi$-flux single-electron hopping amplitude $\kappa=0.4$, pair amplitude $J=1$,
Holstein coupling $g=0.6$ and oscillator cutoff $\fN$, the reflection-positive Gram matrix
of the even graded algebra and of the odd sector have smallest eigenvalues

\begin{center}
\begin{tabular}{@{}rrrr@{}}
\toprule
$\fN$ & $\lambda_{\min}$ (even) & $\lambda_{\min}$ (odd) & $\lambda_{\min}$, trace sector, bond phonon\\
\midrule
4 & $-7.714\cdot10^{-17}$ & $+2.658\cdot10^{-4}$ & $-2.093\cdot10^{-16}$\\
6 & $-2.916\cdot10^{-17}$ & $+2.860\cdot10^{-4}$ & $-3.847\cdot10^{-16}$\\
8 & $-4.167\cdot10^{-17}$ & $+2.868\cdot10^{-4}$ & $-4.095\cdot10^{-16}$\\
10 & $-1.484\cdot10^{-16}$ & $+2.868\cdot10^{-4}$ & $-1.803\cdot10^{-16}$\\
\bottomrule
\end{tabular}
\end{center}

The even-sector and trace-sector values are zero up to the roundoff of the
double-precision Hermitian eigendecomposition, and the odd sector is strictly positive;
the pattern is stable as the cutoff is increased, which is the finite-cutoff
consistency check for Theorem \ref{thm:rp}. This is a check on a two-site system and does
not replace the proof.

\paragraph{Obstructions.}
Two further matrices on the same cluster document the failures of Section \ref{sec:rp}:
taking the oscillator reflection antisymmetrically, $\theta_{\ph}(q_{x})=-q_{\theta(x)}$,
gives $\lambda_{\min}=-0.0263317$ in the odd sector. The failure of part 2 is documented
on the two-site cluster with an oscillator at each site truncated to six levels, at
$\beta=1.5$, $J=0.2$, $U=1$, $\omega=1$ and no Holstein coupling: with a coupling of
strength $1.5$ between the linear oscillator coordinate spanning the reflection plane and
the cross-plane pair bilinear, the smallest eigenvalue of the reflection Gram matrix is
$\lambda_{\min}=-3.943830$.

\paragraph{Certified rows.}
The fifteen rows of Table \ref{tab:rows} are produced by enclosing
$\mathcal{D}(\beta_0)$ in $300\times16$ rectangles and minimizing over them, as described
in Section \ref{sec:finiteT}. The exact lower endpoints of the row enclosures, before the
truncation to five significant digits used in the table, are

\begin{center}
\begin{tabular}{@{}rlll@{}}
\toprule
Row & $m_{0}^{2}\ge$ (exact lower endpoint) & Row & $m_{0}^{2}\ge$ (exact lower endpoint)\\
\midrule
1 & $0.00872325$ & 9 & $0.00671523$\\
2 & $0.0141487$ & 10 & $0.0112208$\\
3 & $0.0502927$ & 11 & $0.000570174$\\
4 & $0.0417303$ & 12 & $0.00614375$\\
5 & $0.0191956$ & 13 & $0.00556189$\\
6 & $0.00439155$ & 14 & $0.00684033$\\
7 & $0.00833005$ & 15 & $0.000558395$\\
8 & $0.0295446$ & & \\
\bottomrule
\end{tabular}
\end{center}

\paragraph{Adversarial modifications.}
Seven semantic modifications of the certified computation are each rejected, in the sense
that the modified computation either produces a negative value where positivity is
required, or fails an exact identity, or violates a variational inequality that the correct
computation satisfies. The seven are: omitting the odd-sublattice phonon parity, which
makes the even-sector Gram matrix negative ($-0.107318$) and leaves the odd sector at the
same negative value as the antisymmetric reflection, $-0.0263317$; dropping the static
phonon term from the sum rule, which violates the exact double-commutator identity by
$0.3729567$ and lets an inadmissible point be certified ($0.0096857$ where the correct
function is $-0.0082233$); reflecting the
oscillators antisymmetrically, which makes the odd-sector Gram matrix negative
($-0.0263317$); reversing the sign of the $\pi$-flux single-electron hopping, as on an
unfrustrated link, which makes the
odd-sector Gram matrix negative ($-0.4392627$); coupling the oscillator coordinate
spanning the reflection plane to the cross-plane pair bilinear, which makes the Gram
matrix negative ($-3.943830$); taking an excessive Holstein coupling $g=2$ at $U=0$,
$\omega=1$, where the certified function is negative ($-1.7159$) while the modified
variational dimer free energy violates the Gibbs bound ($-2.50000$ against an exact value
of $-2.07210$); and pushing the threshold temperature to $T_0=2.5J$, above the
three-dimensional XY ceiling, where the certified function is $-0.63308$ while a reversed
monotonicity in $\beta$ would report $+0.0073385$.

\section{Lattice integrals}\label{app:integrals}

Three Brillouin-zone integrals enter the certified bound. With
$u_{\mathbf p}=\frac13\sum_{\mu}\cos p_{\mu}$, they are
\begin{equation}
\begin{aligned}
I_{W}&:=\int_{[-\pi,\pi]^{3}}\frac{\dd^{3}p}{(2\pi)^{3}}\frac{1}{1+u_{\mathbf p}},\qquad
I_{\sin}:=\int_{[-\pi,\pi]^{3}}\frac{\dd^{3}p}{(2\pi)^{3}}\Big(\sum_{\mu}\sin^{2}p_{\mu}\Big)^{1/2},\\
J_{\mathrm{sq}}(\alpha)&:=\int_{[-\pi,\pi]^{3}}\frac{\dd^{3}p}{(2\pi)^{3}}\frac{(\alpha-u_{\mathbf p})_{+}^{2}}{1+u_{\mathbf p}} .
\end{aligned}
\end{equation}
The first is the cubic Watson integral \cite{watson1939,glasser1977}, for which the
closed form
\begin{equation}
I_{W}=\frac{\sqrt6}{32\pi^{3}}\,\Gamma\Big(\frac{1}{24}\Big)\Gamma\Big(\frac{5}{24}\Big)\Gamma\Big(\frac{7}{24}\Big)\Gamma\Big(\frac{11}{24}\Big)
\label{eq:watson}
\end{equation}
holds after the substitution $\mathbf p\mapsto\mathbf p+\mathbf Q$, which maps
$u_{\mathbf p}$ to $-u_{\mathbf p}$. The enclosing values used in the certified rows are
\begin{equation}
\begin{aligned}
I_{W}&\in[1.516386059,\,1.516386060],\qquad
I_{\sin}\in[1.179837674,\,1.207622426],\\
J_{\mathrm{sq}}\big(\tfrac12\big)&\in[1.408389235,\,1.408950133],\qquad
J_{\mathrm{sq}}\big(\tfrac34\big)\in[2.143651662,\,2.143715091],\\
J_{\mathrm{sq}}(1)&\in[3.065544236,\,3.065544237].
\end{aligned}
\end{equation}
The integrands for $I_{\sin}$ and $J_{\mathrm{sq}}$ are monotone on cells of the
Brillouin zone, which allows them to be enclosed by evaluating cosines and sines at cell
endpoints; the singularity of $J_{\mathrm{sq}}$ at $\mathbf p=\mathbf Q$ is integrable and
is removed by exact subtraction against \eqref{eq:watson}, using
$(\alpha-u)^{2}=(1+\alpha)^{2}-(1+u)\big(2(1+\alpha)-(1+u)\big)$ and the identity that
relates the removed part to $I_{W}$. The Watson integral itself is handled through the
closed form \eqref{eq:watson} with interval arithmetic on the gamma functions.

\sloppy
\bibliographystyle{plain}
\bibliography{refs}

\end{document}
